% locking/locking.tex
% mainfile: ../perfbook.tex
% SPDX-License-Identifier: CC-BY-SA-3.0

\QuickQuizChapter{chp:Locking}{加锁}{qqzlocking}
%
\Epigraph{加锁是最差的通用同步机制，
	  除了所有那些已经被尝试过的
	  其他机制之外。}{向Winston Churchill的记忆
	  以及他所引用的那个人致歉}

在最近的并发研究中，加锁经常扮演反派角色。
\IX{加锁}被指控引发死锁、护送效应、\IX{starvation}、
\IX{unfairness}、\IXpl{data race}，以及各种其他并发罪行。
有趣的是，在生产级共享内存并行软件中，
加锁同样扮演着主力军的角色。
本章将探讨这种反派与英雄之间的对立，
正如
\cref{fig:locking:Locking: Villain or Slob?,%
fig:locking:Locking: Workhorse or Hero?}
中所奇妙地描绘的那样。

这种Jekyll与Hyde式的对立背后有多种原因：

\begin{enumerate}
\item	加锁的许多罪行都有务实的设计解决方案，
	在大多数情况下效果良好，例如：
	\begin{enumerate}
	\item	使用lock hierarchy来避免deadlock。
	\item	Deadlock检测工具，例如Linux内核的
		lockdep工具~\cite{JonathanCorbet2006lockdep}。
	\item	对加锁友好的数据结构，例如
		数组、哈希表和基数树，这些将在
		\cref{chp:Data Structures}中介绍。
	\end{enumerate}
\item	加锁的某些罪行仅在高竞争级别下才成为问题，
	而这种级别只有设计糟糕的程序才会达到。
\item	加锁的某些罪行可以通过将其他同步机制
	与加锁配合使用来避免。
	这些其他机制包括
	统计计数器
	（参见\cref{chp:Counting}），
	引用计数
	（参见\cref{sec:defer:Reference Counting}），
	hazard pointers
	（参见\cref{sec:defer:Hazard Pointers}），
	sequence-locking读者
	（参见\cref{sec:defer:Sequence Locks}），
	RCU
	（参见\cref{sec:defer:Read-Copy Update (RCU)}），
	以及简单的non-blocking数据结构
	（参见\cref{sec:advsync:Non-Blocking Synchronization}）。
\item	直到最近，几乎所有大型共享内存并行程序
	都是专有开发的，因此不容易
	了解到这些务实的解决方案。
\item	加锁对某些软件工件效果极好，
	而对另一些则效果极差。
	曾在加锁效果好的工件上工作过的开发者，
	自然会比那些在加锁效果差的工件上工作过的开发者
	对加锁持更积极的看法，
	这将在\cref{sec:locking:Locking: Hero or Villain?}中讨论。
\item	所有好故事都需要一个反派，而加锁有着作为
	研究论文替罪羊的悠久而光荣的历史。
\end{enumerate}

\QuickQuiz{
	作为替罪羊怎么能被认为是
	光荣的呢？？？
}\QuickQuizAnswer{
	加锁之所以成为研究论文的替罪羊，是因为
	它在实践中被大量使用。
	相比之下，如果没有人使用或关心加锁，大多数研究
	论文甚至不会费心提及它。
}\QuickQuizEnd

本章将概述多种方法，以避免加锁
更严重的罪行。

\begin{figure}
\centering
\resizebox{2in}{!}{\includegraphics{cartoons/r-2014-Locking-the-Slob}}
\caption{加锁：
		  反派还是懒汉？}
\ContributedBy{fig:locking:Locking: Villain or Slob?}{Melissa Broussard}
\end{figure}

\begin{figure}
\centering
\resizebox{2in}{!}{\includegraphics{cartoons/r-2014-Locking-the-Hero}}
\caption{加锁：
		  主力军还是英雄？}
\ContributedBy{fig:locking:Locking: Workhorse or Hero?}{Melissa Broussard}
\end{figure}

\section{保持存活}
\label{sec:locking:Staying Alive}
%
\epigraph{我工作是为了活下去。}{Bette Davis}

鉴于加锁被指控导致死锁和饥饿，
共享内存并行开发者的一个重要关切就是
简单地保持存活。
因此，以下各节将介绍死锁、\IX{livelock}、饥饿、
unfairness和低效率。

\subsection{死锁}
\label{sec:locking:Deadlock}

当一组线程中的每个成员都持有至少一个锁，同时又在等待
该组中另一个成员所持有的锁时，就会发生\IXB{Deadlock}。
即使在只包含单个线程的组中，当该线程试图获取它已经持有的
非递归锁时，这种情况也会发生。
因此，即使只有一个线程和一个锁，也可能发生死锁！

如果没有某种外部干预，死锁将永远持续。
任何线程都无法获取它正在等待的锁，直到持有该锁的线程释放它，
但持有该锁的线程无法释放它，直到该持有线程获取到
它正在等待的锁。

\begin{figure}
\centering
\resizebox{1.5in}{!}{\includegraphics{locking/DeadlockCycle}}
\caption{Deadlock循环}
\label{fig:locking:Deadlock Cycle}
\end{figure}

我们可以用节点表示线程和锁来创建死锁场景的有向图表示，
如\cref{fig:locking:Deadlock Cycle}所示。
从锁指向线程的箭头表示该线程持有该锁，
例如，线程~B持有Lock~2和Lock~4。
从线程指向锁的箭头表示该线程正在等待该锁，
例如，线程~B正在等待Lock~3。

一个死锁场景总是包含至少一个死锁循环。
在\cref{fig:locking:Deadlock Cycle}中，这个循环是
线程~B、Lock~3、线程~C、Lock~4，然后回到线程~B。

\QuickQuiz{
	但是基于锁的死锁的定义只是说每个线程都持有
	至少一个锁并等待另一个被某个线程持有的锁。
	你怎么知道存在一个循环？
}\QuickQuizAnswer{
	假设图中没有循环。
	那么我们将有一个有向无环图（DAG），它至少有一个叶节点。

	如果这个叶节点是一个锁，那么我们就有一个线程在等待
	一个没有被任何线程持有的锁，这与定义矛盾。
	在这种情况下，线程将立即获取该锁。

	另一方面，如果这个叶节点是一个线程，那么我们就有
	一个没有在等待任何lock的线程，同样与定义矛盾。
	在这种情况下，线程要么正在运行，
	要么被阻塞在非lock的某物上。
	在第一种情况下，如果没有无限循环bug，
	线程最终会释放lock。
	在第二种情况下，如果没有唤醒失败的bug，
	线程最终会被唤醒并释放lock。\footnote{
		当然，一种唤醒失败的bug是涉及不仅是lock，
		还有非lock资源的deadlock。
		但问题确实说的是"基于lock的deadlock"！}

	因此，根据基于lock的deadlock的这个定义，
	对应的图中必定存在一个循环。
}\QuickQuizEnd

虽然有一些软件环境（如数据库系统）可以从已有的deadlock中恢复，
但这种方法要求杀死其中一个线程或从其中一个线程强制窃取lock。
这种杀死和强制窃取对事务来说效果很好，
但对于内核和应用层的加锁使用往往有问题：
处理由此产生的部分更新的结构可能极其复杂、危险且容易出错。

因此，内核和应用程序应该避免deadlock。
Deadlock避免策略包括locking hierarchy
（\cref{sec:locking:Locking Hierarchies}）、
局部locking hierarchy
（\cref{sec:locking:Local Locking Hierarchies}）、
分层locking hierarchy
（\cref{sec:locking:Layered Locking Hierarchies}）、
时序locking hierarchy
（\cref{sec:locking:Temporal Locking Hierarchies}）、
处理包含lock指针的API的策略
（\cref{sec:locking:Locking Hierarchies and Pointers to Locks}）、
条件加锁
（\cref{sec:locking:Conditional Locking}）、
先获取所有需要的lock
（\cref{sec:locking:Acquire Needed Locks First}）、
一次一个lock的设计
（\cref{sec:locking:Single-Lock-at-a-Time Designs}）、
以及信号/中断处理程序的策略
（\cref{sec:locking:Signal/Interrupt Handlers}）。
虽然没有一种deadlock避免策略能在所有情况下完美工作，
但有很好的工具可供选择。

\subsubsection{锁层次结构}
\label{sec:locking:Locking Hierarchies}

加锁层次结构对锁进行排序，并禁止以错误的顺序获取锁。
在\cref{fig:locking:Deadlock Cycle}中，
我们可以按数字顺序排列锁，从而禁止线程在已经持有
相同或更高编号的锁时获取给定的锁。
线程~B违反了这个层次结构，因为它试图在持有Lock~4的同时
获取Lock~3。
这种违规导致了死锁的发生。

再次强调，要应用加锁层次结构，需要对锁排序并禁止
乱序的锁获取。
对于不同类型的锁，仔细考虑从一种类型到下一种类型的层次结构
是很有帮助的。
对于同一类型锁的多个实例，例如搜索树中的每节点锁，
传统方法是按要获取的锁的地址顺序进行锁获取。
无论哪种方式，在大型程序中，使用Linux内核
\co{lockdep}~\cite{JonathanCorbet2006lockdep}
等工具来强制执行你的加锁层次结构是明智的。

\subsubsection{局部锁层次结构}
\label{sec:locking:Local Locking Hierarchies}

然而，加锁层次结构的全局性使其难以应用于库函数。
毕竟，当使用给定库函数的程序尚未编写时，
可怜的库函数实现者怎么可能遵循尚未定义的加锁层次结构呢？

一种特殊（但常见）的情况是库函数不调用任何调用者的代码。
在这种情况下，在持有库的任何lock时永远不会获取调用者的lock，
因此不可能存在同时包含库和调用者的lock的deadlock循环。

\QuickQuiz{
	这条规则有什么例外情况吗？即使库代码从不调用
	调用者的任何函数，是否真的可能存在同时包含
	库和调用者的锁的死锁循环？
}\QuickQuizAnswer{
	确实有！
	以下是其中一些：
	\begin{enumerate}
	\item	如果库函数的参数之一是指向锁的指针，
		且该库函数在持有其自身锁时获取调用者的锁，
		那么我们可能会有一个涉及调用者和库锁的
		死锁循环。
	\item	如果库函数之一返回指向锁的指针，
		该锁被调用者获取，且调用者在持有库的锁时
		获取其自身的锁之一，我们又可能有一个涉及
		调用者和库锁的死锁循环。
	\item	如果库函数之一获取了一个锁然后在仍然持有
		该锁的情况下返回，且调用者获取其自身的锁之一，
		我们又有了另一种创建涉及调用者和库锁的
		死锁循环的方式。
	\item	如果调用者有一个获取锁的信号处理程序，
		那么死锁循环可以涉及调用者和库的锁。
		然而在这种情况下，库的锁是死锁循环中的
		无辜旁观者。
		话虽如此，请注意在信号处理程序中获取锁
		在许多环境中是不允许的——这不仅仅是一个坏主意，
		它是不被支持的。
		但如果你确实必须在信号处理程序中获取锁，
		请确保在线程上下文中持有该锁时阻塞该信号，
		以及在持有该锁时获取的任何其他锁时也要阻塞。
	\end{enumerate}
}\QuickQuizEnd

但假设库函数确实调用了调用者的代码。
例如，\co{qsort()}调用调用者提供的比较函数。
通常，这个比较函数将操作不变的局部数据，因此不需要获取锁，
如\cref{fig:locking:No qsort() Compare-Function Locking}所示。
但也许有人疯狂到要对键值正在变化的集合进行排序，
从而要求比较函数获取锁，这可能导致死锁，
如\cref{fig:locking:Without qsort() Local Locking Hierarchy}所示。
库函数如何避免这种死锁？

这种情况下的黄金规则是"在调用未知代码之前释放所有锁。"
要遵循这条规则，\co{qsort()}函数必须在调用比较函数之前
释放其所有锁。
这样\co{qsort()}在比较函数获取调用者的任何锁时
不会持有其任何锁，从而避免了死锁。

\QuickQuiz{
	但如果\co{qsort()}在调用比较函数之前释放了其所有lock，
	它如何防止与其他\co{qsort()}线程的竞争？
}\QuickQuizAnswer{
	通过私有化正在比较的数据元素
	（如\cref{chp:Data Ownership}中所讨论的），
	或者通过使用延迟机制，如引用计数
	（如\cref{chp:Deferred Processing}中所讨论的）。
	或者通过使用分层locking hierarchy，如
	\cref{sec:locking:Layered Locking Hierarchies}中所描述的。

	另一方面，在一个正在排序的列表中更改键值
	至少是相当大胆的。
}\QuickQuizEnd

\begin{figure}
\centering
\resizebox{3in}{!}{\includegraphics{locking/NoLockHierarchy}}
\caption{无\tco{qsort()}比较函数加锁}
\label{fig:locking:No qsort() Compare-Function Locking}
\end{figure}

\begin{figure}
\centering
\resizebox{3in}{!}{\includegraphics{locking/NonLocalLockHierarchy}}
\caption{无\tco{qsort()}局部Locking Hierarchy}
\label{fig:locking:Without qsort() Local Locking Hierarchy}
\end{figure}

\begin{figure}
\centering
\resizebox{3in}{!}{\includegraphics{locking/LocalLockHierarchy}}
\caption{\tco{qsort()}的局部Locking Hierarchy}
\label{fig:locking:Local Locking Hierarchy for qsort()}
\end{figure}

要了解局部locking hierarchy的好处，请比较
\cref{fig:locking:Without qsort() Local Locking Hierarchy,%
fig:locking:Local Locking Hierarchy for qsort()}。
在两个图中，应用程序函数\co{foo()}和\co{bar()}
分别在持有Lock~A和Lock~B时调用\co{qsort()}。
因为这是\co{qsort()}的并行实现，所以它获取了Lock~C\@。
函数\co{foo()}将函数\co{cmp()}传递给\co{qsort()}，
而\co{cmp()}获取Lock~B\@。
函数\co{bar()}将一个简单的整数比较函数（未显示）
传递给\co{qsort()}，这个简单函数不获取任何lock。

现在，如果\co{qsort()}在调用\co{cmp()}时持有Lock~C，违反了
上述释放所有lock的黄金规则，如
\cref{fig:locking:Without qsort() Local Locking Hierarchy}所示，
就可能发生deadlock。
要理解这一点，假设一个线程调用\co{foo()}，同时另一个线程
并发调用\co{bar()}。
第一个线程将获取Lock~A，第二个线程将获取Lock~B\@。
如果第一个线程调用\co{qsort()}获取了Lock~C，那么它在调用
\co{cmp()}时将无法获取Lock~B。
但第一个线程持有Lock~C，所以第二个线程调用\co{qsort()}
将无法获取它，因此也无法释放Lock~B，导致deadlock。

相比之下，如果\co{qsort()}在调用比较函数之前释放Lock~C——
从\co{qsort()}的角度看这是未知代码——
则如\cref{fig:locking:Local Locking Hierarchy for qsort()}所示，
死锁被避免了。

如果每个模块在调用未知代码之前释放所有锁，那么只要每个模块
分别避免死锁，就可以避免死锁。
因此，这条规则大大简化了死锁分析，并大大提高了模块化程度。

然而，这条黄金规则也带有警告。
当你释放那些lock时，它们所保护的任何状态都可能发生任意变化，
而函数的调用者太容易忘记这些变化，导致微妙且难以重现的bug。
因为\co{qsort()}的比较函数很少获取lock，
让我们换一个不同的例子。

考虑\cref{lst:locking:Recursive Tree Iterator}（\path{rec_tree_itr.c}）
中的递归树迭代器。
该迭代器访问树中的每个节点，调用用户的回调函数。
在调用之前释放树lock，在返回后重新获取。
这段代码做了危险的假设：
\begin{enumerate*}[(1)]
\item	当前节点的子节点数量没有变化，
\item	递归存储在栈上的祖先节点仍然存在，以及
\item	被访问的节点本身没有被删除和释放。
\end{enumerate*}
如果一个线程调用\co{tree_add()}而另一个线程释放树的lock
来运行回调函数，就可能遇到其中一些危险。

\QuickQuiz{
	所以迭代线程可能观察到也可能没有观察到添加的子节点。
	这有什么大不了的？
}\QuickQuizAnswer{
	这种情况至少有两个危险。

	一个确实是子节点的数量可能被观察到已更改也可能没有。
	虽然这与\co{tree_add()}在迭代器启动之前或之后被调用
	是一致的，但最好不要将其留给编译器的随意处理。
	一个更严重的问题是\co{realloc()}可能无法原地扩展数组，
	导致堆释放迭代器使用的那个并用另一个内存块替换。
	如果\co{children}指针没有被重新读取，那么迭代线程
	将访问无效内存（要么是空闲的，要么是被回收的）。
}\QuickQuizEnd

一种策略是确保尽管lock被释放，状态仍然被保留，
例如，通过获取节点的引用来防止其被释放。
或者，可以在回调函数返回后重新获取lock时重新初始化状态。

\begin{listing}
\input{CodeSamples/locking/rec_tree_itr@tree_for_each.fcv}
\caption{递归树迭代器}
\label{lst:locking:Recursive Tree Iterator}
\end{listing}

\subsubsection{分层锁层次结构}
\label{sec:locking:Layered Locking Hierarchies}

\begin{figure}
\centering
\resizebox{3in}{!}{\includegraphics{locking/LayeredLockHierarchy}}
\caption{\tco{qsort()}的分层Locking Hierarchy}
\label{fig:locking:Layered Locking Hierarchy for qsort()}
\end{figure}

不幸的是，一方面保留状态可能不可行，
另一方面重新初始化也可能不可行，从而排除了在调用未知代码之前
释放所有锁的局部加锁层次结构。
但是，我们可以代之以构建一个分层加锁层次结构，如
\cref{fig:locking:Layered Locking Hierarchy for qsort()}所示。
这里，\co{cmp()}函数使用一个新的Lock~D，它在所有
Lock~A、B和C之后获取，从而避免了死锁。
因此，我们在全局死锁层次结构中有三层：
第一层包含Lock~A和Lock~B，第二层包含Lock~C，
第三层包含Lock~D\@。

\begin{listing}
\input{CodeSamples/locking/locked_list@start_next.fcv}
\caption{并发列表迭代器}
\label{lst:locking:Concurrent List Iterator}
\end{listing}

请注意，通常无法机械地修改\co{cmp()}以使用新的Lock~D\@。
恰恰相反：通常需要进行深层的设计级修改。
然而，为了避免deadlock而付出这些修改的努力
通常是一个很小的代价。
更重要的是，这种潜在的deadlock最好在设计阶段就被检测到，
在生成任何代码之前！

另一个在调用未知代码之前释放所有lock不切实际的例子是，
设想一个链表的迭代器，如
\cref{lst:locking:Concurrent List Iterator}（\path{locked_list.c}）所示。
\co{list_start()}函数获取列表上的lock并返回第一个元素（如果有的话），
而\co{list_next()}要么返回列表中下一个元素的指针，
要么在到达列表末尾时释放lock并返回\co{NULL}。

\begin{listing}
\input{CodeSamples/locking/locked_list@list_print.fcv}
\caption{并发列表迭代器使用}
\label{lst:locking:Concurrent List Iterator Usage}
\end{listing}

\begin{fcvref}[ln:locking:locked_list:list_print:ints]
\Cref{lst:locking:Concurrent List Iterator Usage}展示了
如何使用这个列表迭代器。
\Clnrefrange{b}{e}定义了包含单个整数的\co{list_ints}元素，
\end{fcvref}
\begin{fcvref}[ln:locking:locked_list:list_print:print]
而\clnrefrange{b}{e}展示了如何遍历列表。
\Clnref{start}对列表加lock并获取指向第一个元素的指针，
\clnref{entry}提供指向我们封装的\co{list_ints}结构的指针，
\clnref{print}打印相应的整数，
\clnref{next}移动到下一个元素。
这相当简单，并隐藏了所有的加锁。
\end{fcvref}

也就是说，只要处理每个列表元素的代码本身不获取
在其他对\co{list_start()}或\co{list_next()}的调用中被持有的lock，
加锁就保持隐藏。否则就会导致deadlock。
我们可以通过分层locking hierarchy来避免deadlock，
将列表迭代器加锁纳入考虑。

这种分层方法可以扩展到任意多的层次，
但每增加一层都会增加加锁设计的复杂性。
这种复杂性的增加对于某些类型的面向对象设计来说特别不方便，
其中控制在一大组对象之间以无纪律的方式来回传递。\footnote{
	这种情况的一个名称是"面向对象的意大利面条代码。"}
面向对象设计的习惯与避免deadlock的需求之间的这种不匹配
是并行编程被某些人认为如此困难的一个重要原因。

一些高度分层locking hierarchy的替代方案在
\cref{chp:Deferred Processing}中介绍。

\subsubsection{时序锁层次结构}
\label{sec:locking:Temporal Locking Hierarchies}

避免deadlock的一种方法是推迟获取冲突lock中的一个。
这种方法在Linux内核的RCU中使用，其\apik{call_rcu()}
函数由Linux内核调度器在持有其lock时调用。
这意味着\co{call_rcu()}不能总是安全地调用调度器
来执行唤醒操作，例如，为了唤醒一个RCU内核线程以启动
由\co{call_rcu()}排队的回调所需的新宽限期。

\QuickQuiz{
	你说"不能总是安全地调用调度器"是什么意思？
	\co{call_rcu()}要么能要么不能安全地调用调度器，
	不是吗？
}\QuickQuizAnswer{
	这确实取决于情况。

	调度器lock总是在中断禁用的情况下持有。
	因此，如果\co{call_rcu()}在中断启用时被调用，
	没有调度器lock被持有，\co{call_rcu()}
	可以安全地调用调度器。
	否则，如果中断被禁用，调度器lock之一\emph{可能}被持有，
	所以\co{call_rcu()}必须谨慎行事，
	避免调用调度器。
}\QuickQuizEnd

然而，宽限期持续许多毫秒，所以在启动新的宽限期之前再等一毫秒
通常不是问题。
因此，如果\co{call_rcu()}检测到与调度器可能存在的死锁，
它会安排稍后启动新的宽限期，要么在定时器处理程序中，
要么在调度器时钟中断处理程序中，具体取决于配置。
因为在任何一个处理程序中都不会持有调度器锁，
死锁被成功避免。

总体方法是通过将锁获取推迟到不持有任何锁的环境中
来遵守加锁层次结构。

\subsubsection{锁层次结构与锁指针}
\label{sec:locking:Locking Hierarchies and Pointers to Locks}

虽然有一些例外，但包含指向lock的指针的外部API
通常是一个设计不良的API\@。
将内部lock交给其他软件组件毕竟是信息隐藏的反面，
而信息隐藏又是一个关键的设计原则。

\QuickQuiz{
	举出一个常见的将lock指针传递给函数的情况。
}\QuickQuizAnswer{
	当然是加锁原语！
}\QuickQuizEnd

一个例外是进行某种实体移交的函数，
其中调用者的lock必须持有直到移交完成，
但lock必须在函数返回之前释放。
一个这样的函数的例子是POSIX的\apipx{pthread_cond_wait()}
函数，其中传递指向\apipx{pthread_mutex_t}的指针
可以防止由于丢失唤醒而导致的挂起。

\QuickQuiz{
	\co{pthread_cond_wait()}先释放mutex然后重新获取它
	这一事实难道不能消除deadlock的可能性吗？
}\QuickQuizAnswer{
	绝对不能！

	考虑一个程序，它按顺序获取\co{mutex_a}，然后
	\co{mutex_b}，然后将\co{mutex_a}
	传递给\co{pthread_cond_wait()}。
	现在，\co{pthread_cond_wait()}将释放\co{mutex_a}，但
	会在返回之前重新获取它。
	如果其他线程在此期间获取了\co{mutex_a}
	然后阻塞在\co{mutex_b}上，程序将deadlock。
}\QuickQuizEnd

简而言之，如果你发现自己正在导出一个以lock指针
作为参数或返回值的API，请帮自己一个忙，
仔细重新考虑你的API设计。
这可能确实是正确的做法，但经验表明这不太可能。

\subsubsection{条件加锁}
\label{sec:locking:Conditional Locking}

但假设不存在合理的locking hierarchy。
这在现实中可能发生，例如，在某些类型的分层网络协议栈中，
数据包在两个方向上流动，例如，在分布式lock管理器的实现中。
在网络的情况下，当将数据包从一层传递到另一层时，
可能需要持有两层的lock。
鉴于数据包在协议栈中上下移动，这是导致deadlock的绝佳配方，
如\cref{lst:locking:Protocol Layering and Deadlock}所示。
\begin{fcvref}[ln:locking:Protocol Layering and Deadlock]
这里，一个向下移动到线路的数据包必须以错误的顺序获取
下一层的lock。
鉴于从线路向上移动的数据包正在按顺序获取lock，
列表中\clnref{acq}处的lock获取可能导致deadlock。
\end{fcvref}

\begin{listing}
\begin{fcvlabel}[ln:locking:Protocol Layering and Deadlock]
\begin{VerbatimL}[commandchars=\\\{\}]
spin_lock(&lock2);
layer_2_processing(pkt);
nextlayer = layer_1(pkt);
spin_lock(&nextlayer->lock1);	\lnlbl{acq}
spin_unlock(&lock2);
layer_1_processing(pkt);
spin_unlock(&nextlayer->lock1);
\end{VerbatimL}
\end{fcvlabel}
\caption{协议分层与Deadlock}
\label{lst:locking:Protocol Layering and Deadlock}
\end{listing}

在这种情况下避免deadlock的一种方法是施加一个locking hierarchy，
但当需要以错误的顺序获取lock时，条件性地获取它，
如\cref{lst:locking:Avoiding Deadlock Via Conditional Locking}所示。
\begin{fcvref}[ln:locking:Avoiding Deadlock Via Conditional Locking]
不是无条件地获取第1层的lock，\clnref{trylock}
使用\apik{spin_trylock()}原语条件性地获取lock。
这个原语在lock可用时立即获取（返回非零值），
否则返回零而不获取lock。

\begin{listing}
\begin{fcvlabel}[ln:locking:Avoiding Deadlock Via Conditional Locking]
\begin{VerbatimL}[commandchars=\\\[\]]
retry:
	spin_lock(&lock2);
	layer_2_processing(pkt);
	nextlayer = layer_1(pkt);
	if (!spin_trylock(&nextlayer->lock1)) {	\lnlbl[trylock]
		spin_unlock(&lock2);		\lnlbl[rel2]
		spin_lock(&nextlayer->lock1);	\lnlbl[acq1]
		spin_lock(&lock2);		\lnlbl[acq2]
		if (layer_1(pkt) != nextlayer) {\lnlbl[recheck]
			spin_unlock(&nextlayer->lock1);
			spin_unlock(&lock2);
			goto retry;
		}
	}
	spin_unlock(&lock2);
	layer_1_processing(pkt);		\lnlbl[l1_proc]
	spin_unlock(&nextlayer->lock1);
\end{VerbatimL}
\end{fcvlabel}
\caption{通过条件加锁避免Deadlock}
\label{lst:locking:Avoiding Deadlock Via Conditional Locking}
\end{listing}

如果\co{spin_trylock()}成功，\clnref{l1_proc}执行所需的
第1层处理。
否则，\clnref{rel2}释放lock，
\clnref{acq1,acq2}按正确顺序获取它们。
不幸的是，系统上可能有多个网络设备
（例如，以太网和WiFi），因此\co{layer_1()}
函数必须做出路由决策。
这个决策可能随时改变，特别是当系统
是移动的时候。\footnote{
	而且，与1900年代不同，移动性是常见情况。}
因此，\clnref{recheck}必须重新检查决策，如果它已更改，
必须释放lock并重新开始。
\end{fcvref}

\QuickQuizSeries{%
\QuickQuizB{
	从\cref{lst:locking:Protocol Layering and Deadlock}到
	\cref{lst:locking:Avoiding Deadlock Via Conditional Locking}
	的转换可以普遍应用吗？
}\QuickQuizAnswerB{
	绝对不能！

	这个转换假设\co{layer_2_processing()}函数是幂等的，
	因为当\co{layer_1()}路由决策改变时，
	它可能在同一个数据包上被执行多次。
	因此，在现实中，这种转换可能变得任意复杂。
}\QuickQuizEndB
%
\QuickQuizE{
	但\cref{lst:locking:Avoiding Deadlock Via Conditional Locking}
	中的复杂性是值得的，因为它避免了死锁，对吧？
}\QuickQuizAnswerE{
	也许。

	如果\co{layer_1()}中的路由决策变化得足够频繁，
	代码将总是重试，永远无法向前推进。
	如果没有线程能向前推进，这被称为"\IX{livelock}"，
	如果某些线程能向前推进但其他线程不能，
	则称为"\IX{starvation}"
	（参见\cref{sec:locking:Livelock and Starvation}）。
}\QuickQuizEndE
}

\subsubsection{先获取所需锁}
\label{sec:locking:Acquire Needed Locks First}

在条件加锁的一个重要特殊情况中，所有需要的lock
在任何处理之前就被获取，其中需要的lock可以通过
对涉及的数据结构的地址进行哈希来识别。
在这种情况下，处理不需要是幂等的：
如果发现在不先释放已获取的lock的情况下无法获取给定的lock，
只需释放所有lock并重试。
只有在持有所有需要的lock之后，才会执行任何处理。

然而，这个过程可能导致\emph{活锁}，这将在
\cref{sec:locking:Livelock and Starvation}中讨论。

\QuickQuiz{
	使用\cref{sec:locking:Acquire Needed Locks First}中描述的
	"先获取所需锁"方法时，如何避免活锁？
}\QuickQuizAnswer{
	提供一个额外的全局锁。
	如果给定线程已经反复尝试并且未能获取所需的锁，
	那么让该线程无条件地获取新的全局锁，
	然后无条件地获取任何所需的锁。
	（由Doug Lea建议。）
}\QuickQuizEnd

一个相关的方法，两阶段加锁~\cite{PhilipABernstein1987}，
在事务数据库系统中已经有了长期的生产使用。
在两阶段加锁事务的第一阶段，锁被获取但不被释放。
一旦所有需要的锁都被获取，事务就进入第二阶段，
锁被释放但不被获取。
这种加锁方法允许数据库为其事务提供可串行化保证，
换句话说，保证事务看到和产生的所有值
与所有事务的某个全局排序一致。
许多这样的系统依赖于中止事务的能力，尽管
通过避免在获取所有需要的锁之前对共享数据进行任何更改
可以简化这一点。
活锁和死锁在这类系统中是问题，但实用的解决方案
可以在许多数据库教科书中找到。

\subsubsection{一次一个锁的设计}
\label{sec:locking:Single-Lock-at-a-Time Designs}

在某些情况下，可以避免嵌套锁，从而避免死锁。
例如，如果一个问题是完全可分区的，可以为每个分区分配一个锁。
然后，在给定分区上工作的线程只需要获取一个对应的锁。
因为没有线程同时持有超过一个锁，死锁是不可能的。

然而，必须有某种机制来确保在没有持有任何锁期间
所需的数据结构继续存在。
一种这样的机制在
\cref{sec:locking:Lock-Based Existence Guarantees}中讨论，
然而，一个重要的限制条件是被处理的问题必须能够被分区。
\Cref{chp:Data Ownership}对数据所有权进行了更详细的讨论，
\cref{chp:Deferred Processing}中还介绍了其他几种方法。

\subsubsection{信号/中断处理程序}
\label{sec:locking:Signal/Interrupt Handlers}

涉及信号处理程序的deadlock通常被快速否定，因为
从信号处理程序内调用\apipx{pthread_mutex_lock()}
是不合法的~\cite{OpenGroup1997pthreads}。
然而，手工构造\emph{可以}从信号处理程序中调用的
加锁原语是可能的（尽管通常不明智）。
此外，几乎所有操作系统内核都允许在中断处理程序中获取锁，
中断处理程序类似于信号处理程序。

诀窍是在获取任何可能在信号（或中断）处理程序中获取的锁时
阻塞信号（或禁用中断，视情况而定）。
此外，如果持有这样的lock，试图获取任何曾在信号处理程序
之外不阻塞信号获取的lock是非法的。

\QuickQuiz{
	假设Lock~A从不在信号处理程序中获取，
	但Lock~B在线程上下文和信号处理程序中都被获取。
	进一步假设Lock~A有时在信号未阻塞的情况下获取。
	为什么在持有Lock~B时获取Lock~A是非法的？
}\QuickQuizAnswer{
	因为这会导致deadlock。
	鉴于Lock~A有时在信号处理程序之外不阻塞信号的情况下被持有，
	在持有此lock时可能会处理一个信号。
	相应的信号处理程序可能然后获取Lock~B，
	因此Lock~B在持有Lock~A的情况下被获取。
	因此，如果我们在持有Lock~B时也获取Lock~A，
	我们将有一个deadlock循环。
	注意，即使在持有Lock~B时信号被阻塞，
	这个问题仍然存在。

	这是在中断或信号处理程序中获取lock时要非常小心的
	又一个原因。
	但Linux内核的lock依赖检查器了解这种情况
	以及许多其他情况，所以请充分利用它！
}\QuickQuizEnd

如果一个lock被多个信号的处理程序获取，那么每当获取该lock时，
必须阻塞这些信号中的每一个，
即使该lock是在信号处理程序内获取的。

\QuickQuiz{
	如何在信号处理程序中合法地阻塞信号？
}\QuickQuizAnswer{
	最简单和最快的方法之一是使用
	传递给\co{sigaction()}设置信号时
	\co{struct sigaction}的\co{sa_mask}字段。
}\QuickQuizEnd

不幸的是，在某些操作系统中阻塞和解除阻塞信号可能很昂贵，
特别是包括Linux在内，因此性能方面的考虑通常意味着
在信号处理程序中获取的lock只在信号处理程序中获取，
并且使用无锁同步机制在应用程序代码和信号处理程序之间通信。

或者完全避免信号处理程序，除非处理致命错误。

\QuickQuiz{
	如果在信号处理程序中获取lock是这么糟糕的主意，
	为什么还要讨论使其安全的方法？
}\QuickQuizAnswer{
	因为同样的规则适用于操作系统内核和某些嵌入式应用程序中
	使用的中断处理程序。

	在许多应用程序环境中，在信号处理程序中获取lock
	是不被赞成的~\cite{OpenGroup1997pthreads}。
	然而，这并不能阻止聪明的开发者（也许不明智地）
	使用原子操作构造自制lock。
	而原子操作在许多情况下在信号处理程序中是完全合法的。
}\QuickQuizEnd

\subsubsection{讨论}
\label{sec:locking:Locking Hierarchy Discussion}

共享内存并行程序员有大量的deadlock避免策略可用，但有些
顺序程序，没有一种策略能很好地适用。
这是专家程序员的工具箱中有不止一种工具的原因之一：
加锁是强大的并发工具，但有些工作用其他工具更好解决。

\QuickQuiz{
	给定一个面向对象的应用程序，它在一组对象之间自由传递控制，
	使得没有直接的locking hierarchy，\footnote{
		也被称为"面向对象的意大利面条代码。"}
	无论是分层的还是其他的，这个应用程序如何被并行化？
}\QuickQuizAnswer{
	有多种方法：
	\begin{enumerate}
	\item	在通过模拟进行参数搜索的情况下，
		将运行大量模拟以收敛到（例如）机械或电气设备的
		良好设计，让模拟保持单线程，但并行运行模拟的
		许多实例。
		这保留了面向对象的设计，并在更高层次获得
		并行性，且可能也避免了死锁和同步开销。
	\item	将对象分成组，使得在给定时间不需要操作
		多个组中的对象。
		然后将一个锁与每个组关联。
		这是一次一个锁设计的例子，在
		\cref{sec:locking:Single-Lock-at-a-Time Designs}中讨论。
	\item	将对象分成组，使线程都可以按某种组间顺序
		操作组中的对象。
		然后将一个锁与每个组关联，并对组施加
		加锁层次结构。
	\item	对锁施加任意选择的层次结构，
		然后在需要以错误顺序获取锁时使用条件加锁，
		如\cref{sec:locking:Conditional Locking}中所讨论的。
	\item	在执行一组给定的操作之前，预测将获取哪些锁，
		并在实际执行任何更新之前尝试获取它们。
		如果预测结果不正确，放弃所有锁并
		使用包含经验教训的更新预测重试。
		这种方法在
		\cref{sec:locking:Acquire Needed Locks First}中讨论。
	\item	使用事务内存。
		这种方法有许多优点和缺点，
		将在\crefrange{sec:future:Transactional Memory}{sec:future:Hardware Transactional Memory}中讨论。
	\item	重构应用程序使其更适合并发。
		这可能也有使应用程序即使在单线程下运行更快的
		副作用，但也可能使修改应用程序变得更困难。
	\item	除了加锁之外，还使用后面章节中的技术。
	\end{enumerate}
}\QuickQuizEnd

然而，本节描述的策略已被证明在许多场景中相当有用。

\subsection{活锁与饥饿}\ucindex{Livelock|BF}\ucindex{Starvation|BF}
\label{sec:locking:Livelock and Starvation}

虽然条件加锁可以是一种有效的死锁避免机制，但它可能被滥用。
例如考虑\cref{lst:locking:Abusing Conditional Locking}中
展示的美丽对称的例子。
这个例子的美丽隐藏了一个丑陋的活锁。
要理解这一点，考虑以下事件序列：

\begin{listing}
\begin{fcvlabel}[ln:locking:Abusing Conditional Locking]
\begin{VerbatimL}[commandchars=\\\[\]]
void thread1(void)
{
retry:					\lnlbl[thr1:retry]
	spin_lock(&lock1);		\lnlbl[thr1:acq1]
	do_one_thing();
	if (!spin_trylock(&lock2)) {	\lnlbl[thr1:try2]
		spin_unlock(&lock1);    \lnlbl[thr1:rel1]
		goto retry;
	}
	do_another_thing();
	spin_unlock(&lock2);
	spin_unlock(&lock1);
}

void thread2(void)
{
retry:					\lnlbl[thr2:retry]
	spin_lock(&lock2);		\lnlbl[thr2:acq2]
	do_a_third_thing();
	if (!spin_trylock(&lock1)) {	\lnlbl[thr2:try1]
		spin_unlock(&lock2);	\lnlbl[thr2:rel2]
		goto retry;
	}
	do_a_fourth_thing();
	spin_unlock(&lock1);
	spin_unlock(&lock2);
}
\end{VerbatimL}
\end{fcvlabel}
\caption{滥用条件加锁}
\label{lst:locking:Abusing Conditional Locking}
\end{listing}

\begin{fcvref}[ln:locking:Abusing Conditional Locking]
\begin{enumerate}
\item	线程~1在\clnref{thr1:acq1}获取\co{lock1}，然后调用
	\co{do_one_thing()}。
\item	线程~2在\clnref{thr2:acq2}获取\co{lock2}，然后调用
	\co{do_a_third_thing()}。
\item	线程~1试图在\clnref{thr1:try2}获取\co{lock2}，
	但因为线程~2持有它而失败。
\item	线程~2试图在\clnref{thr2:try1}获取\co{lock1}，
	但因为线程~1持有它而失败。
\item	线程~1在\clnref{thr1:rel1}释放\co{lock1}，
	然后跳转到\clnref{thr1:retry}处的\co{retry}。
\item	线程~2在\clnref{thr2:rel2}释放\co{lock2}，
	并跳转到\clnref{thr2:retry}处的\co{retry}。
\item	活锁之舞从头重复。
\end{enumerate}
\end{fcvref}

\QuickQuiz{
	如何避免\cref{lst:locking:Abusing Conditional Locking}中
	展示的活锁？
}\QuickQuizAnswer{
	\Cref{lst:locking:Avoiding Deadlock Via Conditional Locking}
	提供了一些很好的提示。
	在许多情况下，活锁是一个提示，说明你应该重新审视
	你的加锁设计。
	或者如果你的加锁设计是"自然成长"的，
	就首先审视它。

	话虽如此，Doug Lea提出的一种好而充分的方法是
	使用\cref{sec:locking:Conditional Locking}中描述的条件加锁，
	但将其与在修改共享数据之前先获取所有需要的锁结合起来，
	如\cref{sec:locking:Acquire Needed Locks First}中所描述的。
	如果给定的critical section重试次数过多，
	无条件地获取一个全局锁，然后无条件地获取
	所有需要的锁。
	这既避免了死锁又避免了活锁，并且在假设全局锁
	不需要被太频繁获取的情况下具有合理的可扩展性。
}\QuickQuizEnd

活锁可以被认为是饥饿的一种极端形式，
其中是一组线程在饿死，而不仅仅是其中一个。\footnote{
	不要过于纠结于活锁、饥饿和unfairness
	等术语的精确定义。
	任何导致一组线程无法取得足够前进进度的情况
	都是需要修复的bug，而争论名称不能修复bug。}

\begin{listing}
\begin{fcvlabel}[ln:locking:Conditional Locking and Exponential Backoff]
\begin{VerbatimL}[commandchars=\\\[\]]
void thread1(void)
{
	unsigned int wait = 1;
retry:
	spin_lock(&lock1);
	do_one_thing();
	if (!spin_trylock(&lock2)) {
		spin_unlock(&lock1);
		sleep(wait);
		wait = wait << 1;
		goto retry;
	}
	do_another_thing();
	spin_unlock(&lock2);
	spin_unlock(&lock1);
}

void thread2(void)
{
	unsigned int wait = 1;
retry:
	spin_lock(&lock2);
	do_a_third_thing();
	if (!spin_trylock(&lock1)) {
		spin_unlock(&lock2);
		sleep(wait);
		wait = wait << 1;
		goto retry;
	}
	do_a_fourth_thing();
	spin_unlock(&lock1);
	spin_unlock(&lock2);
}
\end{VerbatimL}
\end{fcvlabel}
\caption{条件加锁与指数退避}
\label{lst:locking:Conditional Locking and Exponential Backoff}
\end{listing}

活锁和饥饿是软件事务内存实现中的严重问题，
因此引入了\emph{竞争管理器}的概念来封装这些问题。
在加锁的情况下，简单的指数退避通常可以解决活锁和饥饿。
其思想是在每次重试之前引入指数增长的延迟，
如\cref{lst:locking:Conditional Locking and Exponential Backoff}所示。

\QuickQuiz{
	你能在\cref{lst:locking:Conditional Locking and Exponential Backoff}
	中的代码里发现什么问题？
}\QuickQuizAnswer{
	以下是几个：
	\begin{enumerate}
	\item	对大多数用途来说，一秒的等待时间太长了。
		等待间隔应该从大约执行critical section
		所需的时间开始，通常在微秒或毫秒范围内。
	\item	代码没有检查溢出。
		另一方面，这个bug被前一个bug抵消了：
		32位的秒数超过50年。
	\end{enumerate}
}\QuickQuizEnd

为了获得更好的结果，退避应该有界限，
而在高竞争情况下，通过排队加锁~\cite{Anderson90}
可以获得更好的结果，这在
\cref{sec:locking:Other Exclusive-Locking Implementations}中进一步讨论。
当然，最好的方法是使用好的并行设计，
通过保持低\IX{lock contention}来避免这些问题。

\subsection{不公平性}
\label{sec:locking:Unfairness}

\begin{figure}
\centering
\resizebox{3in}{!}{\includegraphics{cpu/SystemArch}}
\caption{系统架构与Lock Unfairness}
\label{fig:locking:System Architecture and Lock Unfairness}
\end{figure}

\IXB{Unfairness}可以被认为是饥饿的一种较轻形式，
其中竞争给定锁的线程子集获得了大部分的锁获取机会。
这可能发生在具有共享缓存或\IXacr{numa}特征的机器上，
例如，如\cref{fig:locking:System Architecture and Lock Unfairness}所示。
如果CPU~0释放了一个所有其他CPU都试图获取的锁，
CPU~0和CPU~1之间共享的互连意味着
CPU~1比CPU~2--7有优势。
因此CPU~1可能会获取该锁。
如果CPU~1持有锁的时间足够长，使得CPU~0在CPU~1释放它时
正在请求该锁，反之亦然，锁可以在CPU~0和CPU~1之间穿梭，
绕过CPU~2--7。

\QuickQuiz{
	使用好的并行设计使锁竞争足够低以避免unfairness
	不是更好吗？
}\QuickQuizAnswer{
	在某种意义上是更好的，但有些情况下使用有时会导致
	高锁竞争的设计是合适的。

	例如，设想一个受罕见错误条件影响的系统。
	对于罕见错误条件的持续时间，拥有一个性能和可扩展性
	较差的简单错误处理设计可能比拥有一个只在那些
	罕见错误条件生效时才有用的复杂且难以调试的设计更好。

	话虽如此，通常值得投入一些精力来尝试产生一个
	在错误条件下既简单又高效的设计，例如通过分区问题。
}\QuickQuizEnd

\subsection{低效率}
\label{sec:locking:Inefficiency}

锁是使用原子指令和\IXpl{memory barrier}实现的，
而且经常涉及缓存未命中。
正如我们在\cref{chp:Hardware and its Habits}中看到的，
这些指令非常昂贵，大约比简单指令高两个数量级的开销。
这对加锁来说可能是一个严重的问题：
如果你用锁保护单条指令，你将把开销增加一百倍。
即使假设完美的可扩展性，也需要\emph{一百}个CPU才能跟上
一个不使用加锁执行相同代码的单个CPU。

\begin{figure}
\centering
\resizebox{2.0in}{!}{\includegraphics{locking/sawkerf}}
\caption{锯口}
\label{fig:locking:Saw Kerf}
\end{figure}

这种情况不仅限于加锁。
\Cref{fig:locking:Saw Kerf}
展示了同样的原理如何适用于古老的锯木活动。
如图所示，锯一块板会将板的一小部分（锯片的宽度）转化为锯末。
当然，锁是对时间进行分区而不是锯木头，\footnote{
	也就是说，加锁是时间同步。
	在时间和空间上都进行同步的机制在
	\cref{chp:Deferred Processing}中描述。}
但就像锯木头一样，使用锁对时间进行分区会由于锁开销
和（更糟糕的）锁竞争而浪费一些时间。
一个重要的区别是，如果有人把板锯成太小的碎片，
大部分板变成锯末的结果将立即显而易见。
相比之下，给定的锁获取是否浪费了过多的时间并不总是显而易见的。

这种情况强调了\cref{sec:SMPdesign:Synchronization Granularity}中
讨论的同步粒度权衡的重要性，
特别是\cref{fig:SMPdesign:Synchronization Efficiency}：
粒度太粗会限制可扩展性，而粒度太细则会导致过多的同步开销。

获取锁可能很昂贵，但一旦持有，CPU的缓存是一个有效的
性能助推器，至少对于大型\IXpl{critical section}是这样。
此外，一旦持有锁，被该锁保护的数据可以被锁持有者
访问而不受其他线程的干扰。

\QuickQuiz{
	锁持有者如何被干扰？
}\QuickQuizAnswer{
	如果被锁保护的数据与锁本身在同一缓存行中，
	那么其他CPU获取锁的尝试会导致持有锁的CPU
	产生昂贵的缓存未命中。
	这是\IX{false sharing}的一个特殊情况，当被不同锁保护的
	一对变量碰巧共享一个缓存行时也会发生。
	相比之下，如果锁与它保护的数据在不同的缓存行中，
	持有锁的CPU通常只在第一次访问给定变量时
	才会产生缓存未命中。

	当然，将锁和数据放在不同缓存行的缺点是
	在无竞争情况下代码将产生两次缓存未命中而不是一次。
	一如既往，明智地选择！
}\QuickQuizEnd

Rust编程语言将锁/数据关联更进一步，
允许开发者在锁和它保护的数据之间建立编译器可见的关联~\cite{RalfJung2021RustSafeSysProg}。
当建立了这样的关联后，在没有相应锁保护的情况下尝试访问数据
将产生编译时诊断。
希望这将大大减少这类bug的频率。
当然，这种方法不能直接适用于被加锁的数据分布在
某个数据结构的节点中的情况，或者被加锁的是纯抽象的东西时，
例如，当一小部分状态机转换要被给定锁保护时。
出于这个原因，Rust允许锁与类型而不是数据项关联，
甚至可以不与任何东西关联。
最后一个选项允许Rust模拟传统的加锁使用场景，
但在Rust开发者中不太流行。
也许Rust社区将针对其他加锁使用场景提出其他机制。

\section{锁的类型}
\label{sec:locking:Types of Locks}
%
\epigraph{生活中唯一的锁是你以为你知道但实际不知道的东西。
	  接受你的无知并尝试新事物。}
	 {Dennis Vickers}

锁的类型之多令人惊讶，远超本短章所能详尽的。
以下各节讨论
互斥锁（\cref{sec:locking:Exclusive Locks}）、
读写锁（\cref{sec:locking:Reader-Writer Locks}）、
多角色锁（\cref{sec:locking:Beyond Reader-Writer Locks}）、
以及作用域加锁（\cref{sec:locking:Scoped Locking}）。

\subsection{互斥锁}
\label{sec:locking:Exclusive Locks}

\IXhpl{Exclusive}{lock}顾名思义：
同一时间只有一个线程可以持有该锁。
因此，这样的锁的持有者对该锁保护的所有数据拥有独占访问权，
故得此名。

当然，这一切都假设在访问据称被锁保护的数据时
该锁都被持有。
虽然有一些工具可以帮助
（例如参见\cref{sec:formal:Axiomatic Approaches and Locking}），
确保始终在需要时获取锁的最终责任在开发者身上。

\QuickQuiz{
	让一个互斥锁获取紧接着释放同一个锁，
	即空的critical section，有意义吗？
}\QuickQuizAnswer{
	基于锁的空critical section很少使用，但确实有其用途。
	要点是互斥锁的语义有两个组成部分：
	\begin{enumerate*}[(1)]
	\item 熟悉的数据保护语义和
	\item 消息语义，释放给定锁会通知等待获取同一锁的操作。
	\end{enumerate*}
	空的critical section使用消息组件而不使用数据保护组件。

	这个答案的其余部分提供了一些空critical section
	的用法示例，但是，这些示例应被视为
	``灰色魔法。''\footnote{
		Thanks to Alexey Roytman for this description.}
	因此，空的critical section在实践中几乎从不使用。
	尽管如此，让我们继续深入这个灰色地带\ldots

	空critical section的一个历史用例出现在
	2.4版Linux内核的网络栈中，通过使用一个
	名为\co{brlock}的读端可扩展的读写锁，
	即``大读者锁''。
	这个用例是一种近似read-copy
	update（RCU）语义的方式，RCU在
	\cref{sec:defer:Read-Copy Update (RCU)}中讨论。
	实际上这个Linux内核用例已被
	RCU\@取代。

	空锁critical section习语也可以在某些情况下
	用于减少锁竞争。
	例如，考虑一个多线程用户空间应用程序，其中
	每个线程处理维护在每线程列表中的工作单元，
	线程被禁止触碰彼此的
	列表~\cite{PaulEMcKenney2012EmptyLocks}。
	也可能存在需要所有先前调度的工作单元
	完成后才能进行的更新。
	处理这种情况的一种方法是在每个线程上
	调度一个工作单元，这样当所有这些工作单元
	完成时，更新就可以继续。

	在某些应用程序中，线程可以来去自如。
	例如，每个线程可能对应应用程序的一个用户，
	因此当该用户注销或以其他方式断开连接时
	被移除。
	在许多应用程序中，线程不能原子地离开：
	它们必须使用特定的操作序列
	从应用程序的各个部分显式地解除自身。
	一个特定操作是拒绝接受来自其他线程的更多请求，
	另一个特定操作是处理其列表上
	剩余的工作单元，例如，将这些工作单元
	放入全局工作项处置列表中，
	由剩余线程之一取走。
	（为什么不通过执行每个项目来排空线程的
	工作项列表？
	因为给定的工作项可能会生成更多工作项，
	所以列表无法及时排空。）

	如果应用程序要表现良好并具有良好的可扩展性，
	则需要良好的加锁设计。
	一个常见的解决方案是使用一个全局锁（称之为\co{G}）
	保护整个离开过程
	（也许还有其他事情），
	使用更细粒度的锁保护
	各个解除操作。

	现在，离开的线程必须在处理其列表上的工作之前
	明确拒绝接受更多请求，因为否则
	在处理操作之后可能会有额外的工作到达，
	这将使该处理操作无效。
	因此离开线程的简化伪代码可能如下：

	\begin{enumerate}
	\item	获取锁\co{G}。
	\item	获取保护通信的锁。
	\item	拒绝来自其他线程的进一步通信。
	\item	释放保护通信的锁。
	\item	获取保护全局工作项处置列表的锁。
	\item	将所有待处理的工作项移到全局
		工作项处置列表。
	\item	释放保护全局工作项处置列表的锁。
	\item	释放锁\co{G}。
	\end{enumerate}

	当然，需要等待所有预先存在的工作项的
	线程需要考虑正在离开的线程。
	为了说明这一点，假设该线程在一个正在离开的线程
	拒绝了来自其他线程的进一步通信后
	开始等待所有预先存在的工作项。
	这个线程如何等待离开线程的工作项完成，
	请记住线程不允许
	访问彼此的工作项列表？

	一种直接的方法是让该线程获取\co{G}
	然后获取保护全局工作项处置列表的lock，然后
	将工作项移到自己的列表。
	然后该线程释放两个lock，
	在自己列表的末尾放置一个工作项，
	然后等待它放置在每个线程列表上的
	所有工作项（包括自己的）完成。

	这种方法在许多情况下确实效果很好，但如果
	从全局工作项处置列表中取出每个工作项时
	需要特殊处理，结果可能是
	\co{G}上的过度竞争。
	避免该竞争的一种方法是获取\co{G}然后
	立即释放它。
	然后等待所有先前工作项的过程
	看起来像下面这样：

	\begin{enumerate}
	\item	将全局计数器设为1并将条件
		变量初始化为零。
	\item	向所有线程发送消息，使它们原子地
		增加全局计数器，然后入队一个
		工作项。
		该工作项将原子地递减全局
		计数器，如果结果为零，它将把
		条件变量设为1。
	\item	获取\co{G}，这将等待任何当前正在离开的
		线程完成离开。
		因为一次只有一个线程可以离开，所有
		剩余的线程将已经收到了前一步
		发送的消息。
	\item	释放\co{G}。
	\item	获取保护全局工作项处置列表的lock。
	\item	将所有工作项从全局工作项处置列表
		移到该线程的列表，沿途根据需要处理它们。
	\item	释放保护全局工作项处置列表的lock。
	\item	在该线程的列表上入队一个额外的工作项。
		（和之前一样，这个工作项将原子地递减
		全局计数器，如果结果为零，它将
		把条件变量设为1。）
	\item	等待条件变量变为值1。
	\end{enumerate}

	一旦这个过程完成，所有预先存在的工作项
	保证已经完成。
	空的critical section使用加锁进行消息传递
	以及数据保护。
}\QuickQuizEnd

\QuickQuizLabel{\QlockingQemptycriticalsection}

重要的是要注意，无条件获取互斥锁有两个效果：
\begin{enumerate*}[(1)]
\item 等待该锁的所有先前持有者释放它，以及
\item 阻止任何其他获取尝试，直到锁被释放。
\end{enumerate*}
因此，在锁获取时，该锁的任何并发获取
必须被划分为先前持有者和后续持有者。
不同类型的互斥锁使用不同的划分
策略~\cite{BjoernBrandenburgPhD,Guerraoui:2019:LPA:3319851.3301501}，
例如：

\begin{enumerate}
\item	严格FIFO，较早开始获取的先获取锁。
\item	近似FIFO，足够早开始获取的先获取锁。
\item	优先级级别内的FIFO，高优先级线程比大约同时尝试获取
	锁的任何低优先级线程更早获取锁，但对于相同优先级的
	线程应用某种FIFO排序。
\item	随机，新的锁持有者从所有尝试获取的线程中
	随机选择，不考虑时序。
\item
	不公平，给定的获取可能永远无法获取锁
	（参见\cref{sec:locking:Unfairness}）。
\end{enumerate}

不幸的是，具有更强保证的加锁实现通常会产生更高的开销，
这推动了生产中使用的各种加锁实现。
例如，实时系统通常需要优先级级别内的某种程度的FIFO排序，
以及许多其他方面
（参见\cref{sec:advsync:Event-Driven Real-Time Support}），
而受到高竞争的非实时系统可能只需要足够的排序来避免饥饿，
最后，设计为避免竞争的非实时系统可能根本不需要公平性。

\subsection{读写锁}
\label{sec:locking:Reader-Writer Locks}

\IXhpl{Reader-writer}{lock}~\cite{Courtois71}
允许任意数量的读者同时持有锁，
或者允许单个写者持有锁。
因此，从理论上讲，读写锁应该能为经常读取
而很少写入的数据提供出色的可扩展性。
在实践中，可扩展性将取决于读写锁的实现。

经典的读写锁实现涉及一组原子操作的
计数器和标志。
这种类型的实现与短critical section的互斥锁有同样的问题：
获取和释放锁的开销大约比简单指令的开销
高两个数量级。
当然，如果critical section足够长，获取和释放锁的开销
就变得可以忽略不计。
然而，因为同一时间只有一个线程可以操作锁，
所需的critical section大小会随CPU数量增加。

可以通过使用每线程互斥锁来设计一种对读者更有利的
读写锁~\cite{WilsonCHsieh92a}。
要读取，线程只获取自己的锁。
要写入，线程获取所有锁。
在没有写者的情况下，每个读者只产生原子指令和memory barrier开销，
没有缓存未命中，这对于加锁原语来说相当好。
不幸的是，写者必须产生缓存未命中以及原子指令和
memory barrier开销——乘以线程数量。

简而言之，读写锁在许多情况下非常有用，
但每种类型的实现都有其缺点。
读写锁的典型使用场景涉及非常长的
\IXhpl{read-side}{critical section}，最好以数百微秒
甚至毫秒来衡量。

与互斥锁一样，读写锁的获取在
该锁的所有先前冲突持有者释放它之前无法完成。
如果锁被读持有，则读获取可以立即完成，
但写获取必须等到不再有任何读者持有该锁。
如果锁被写持有，则所有获取都必须等到写者释放锁。
同样与互斥锁一样，不同的读写锁实现
分别为读者和写者提供不同程度的FIFO排序。

但假设大量读者持有锁而一个写者正在等待获取锁。
是否应该允许读者继续获取锁，可能导致写者饥饿？
类似地，假设一个写者持有锁而大量读者和写者
都在等待获取锁。
当当前写者释放锁时，应该将其给读者还是另一个写者？
如果给了读者，在允许下一个写者获取之前
应该允许多少读者获取锁？

这些问题有许多可能的答案，具有不同的复杂性、开销和公平性水平。
不同的实现可能有不同的成本，例如，
某些类型的读写锁在从读持有者切换到写持有者模式时
产生极大的延迟。
以下是一些可能的方法：

\begin{enumerate}
\item
	读者优先实现无条件地优先读者而非写者，
	可能允许写获取被无限期阻塞。
\item	批量公平实现确保当读者和写者都在获取lock时，
	两者都通过批处理获得合理的访问。
	例如，锁可能允许每个CPU五个读者，然后两个写者，
	然后每个CPU再五个读者，依此类推。
\item	写者优先实现无条件地优先写者而非读者，
	可能允许读获取被无限期阻塞。
\end{enumerate}

当然，这些区别只在高lock contention条件下才有意义。

请牢记lock的等待/阻塞双重本质。
这将在\cref{chp:Deferred Processing}的讨论中重新审视，
其中讨论了加锁的可扩展高性能专用替代方案。

\subsection{超越读写锁}
\label{sec:locking:Beyond Reader-Writer Locks}

\begin{table}
\renewcommand*{\arraystretch}{1.2}
\newcommand{\x}{\textcolor{gray!20}{\rule{7pt}{7pt}}}
\newcommand{\rothead}[1]{\begin{picture}(6,65)(0,0)\rotatebox{90}{#1}\end{picture}}
\small
\centering
\begin{tabular}{lcccccc}
	\toprule
	& \rothead{Null (Not Held)}
	& \rothead{Concurrent Read}
	& \rothead{Concurrent Write}
	& \rothead{Protected Read}
	& \rothead{Protected Write}
	& \rothead{Exclusive}
	\\
%				  NL   CR   CW     PR   PW   EX
	\cmidrule(r){1-1} \cmidrule{2-7}
	Null (Not Held)		& \x & \x & \x   & \x & \x & \x \\
	Concurrent Read		& \x & \x & \x   & \x & \x &  X \\
	Concurrent Write	& \x & \x & \x   &  X &  X &  X \\
	Protected Read		& \x & \x &  X   & \x &  X &  X \\
	Protected Write		& \x & \x &  X   &  X &  X &  X \\
	Exclusive		& \x &  X &  X   &  X &  X &  X \\
	\bottomrule
\end{tabular}
\caption{VAX/VMS分布式Lock管理器策略}
\label{tab:locking:VAX/VMS Distributed Lock Manager Policy}
\end{table}

读写锁和互斥锁在其准入策略上有所不同：
互斥锁最多允许一个持有者，而读写锁
允许任意数量的读持有者（但只允许一个写持有者）。
可能的准入策略有非常多种，其中之一是VAX/VMS分布式锁
管理器（DLM）~\cite{Snaman87}的策略，如
\cref{tab:locking:VAX/VMS Distributed Lock Manager Policy}所示。
空白单元格表示兼容模式，而包含"X"的单元格表示不兼容模式。

VAX/VMS DLM使用六种模式。
为了进行比较，互斥锁使用两种模式（未持有和已持有），
而读写锁使用三种模式（未持有、读持有和写持有）。

第一种模式是null，即未持有。
这种模式与所有其他模式兼容，这是预期的：
如果一个线程没有持有锁，它不应该阻止其他线程获取该锁。

第二种模式是concurrent read，它与除exclusive之外的所有其他模式兼容。
concurrent-read模式可能用于在允许更新并发进行的同时
累积数据结构的近似统计。

第三种模式是concurrent write，它与null、
concurrent read和concurrent write兼容。
concurrent-write模式可能用于更新近似统计，
同时仍允许读取和并发更新并发进行。

第四种模式是protected read，它与null、
concurrent read和protected read兼容。
protected-read模式可能用于获取数据结构的一致快照，
同时允许读取但不允许更新并发进行。

第五种模式是protected write，它与null和
concurrent read兼容。
protected-write模式可能用于对数据结构进行更新，
这些更新可能干扰protected reader但可以被concurrent reader容忍。

第六种也是最后一种模式是exclusive，它只与null兼容。
exclusive模式在需要排除所有其他访问时使用。

有趣的是，互斥锁和读写锁
可以由VAX/VMS DLM模拟。
互斥锁只使用null和exclusive模式，而
读写锁可能使用null、protected-read和
protected-write模式。

\QuickQuiz{
	VAX/VMS DLM还有其他方式来模拟
	读写锁吗？
}\QuickQuizAnswer{
	实际上有好几种。
	一种方式是使用null、protected-read和exclusive
	模式。
	另一种方式是使用null、protected-read和
	concurrent-write模式。
	第三种方式是使用null、concurrent-read和
	exclusive模式。
}\QuickQuizEnd

虽然VAX/VMS DLM策略已在分布式数据库中得到
广泛的生产使用，但似乎在共享内存应用程序中
使用不多。
一个可能的原因是分布式数据库较大的通信开销
可以隐藏VAX/VMS DLM更复杂的
准入策略的较大开销。

尽管如此，VAX/VMS DLM是一个有趣的例证，说明
加锁背后的概念可以有多灵活。
它也是对现代DBMS使用的加锁方案的
一个非常简单的介绍，现代DBMS可以有超过三十种加锁模式，
而VAX/VMS只有六种。

\subsection{作用域加锁}
\label{sec:locking:Scoped Locking}

到目前为止讨论的加锁原语需要显式的获取和释放原语，
例如分别是\co{spin_lock()}和\co{spin_unlock()}。
另一种方法是使用面向对象的\IXacrmfst{raii}
模式~\cite{MargaretAEllis1990Cplusplus}。\footnote{
	在\url{https://www.stroustrup.com/bs_faq2.html\#finally}
	中有更清晰的表达。}
这种模式通常应用于C++等语言中的auto变量，
其中相应的\emph{构造函数}在进入
对象的作用域时被调用，相应的\emph{析构函数}在
退出该作用域时被调用。
这可以应用于加锁，让构造函数获取锁，
析构函数释放它。

这种方法可以非常有用，实际上在1990年我确信它
是唯一需要的加锁类型。\footnote{
	我后来在Sequent Computer Systems做并行工作
	很快就让我摒弃了这个错误的观念。}
RAII加锁的一个非常好的特性是你不需要在
退出该作用域的每条代码路径上仔细释放锁，
这个特性可以消除一组麻烦的bug。

然而，RAII加锁也有其黑暗面。
RAII使得封装锁获取和释放变得相当困难，
例如在迭代器中。
在许多迭代器实现中，你希望在
迭代器的``start''函数中获取锁，在迭代器的``stop''
函数中释放它。
RAII加锁要求锁获取和释放必须在
同一层级的作用域中进行，这使得这种封装困难甚至
不可能。

严格的RAII加锁还禁止重叠的critical section，因为
作用域必须嵌套。
这种禁止使得表达许多有用的构造变得困难或不可能，
例如，在多个并发尝试断言事件之间
进行调解的加锁树。
在任意大的一组并发尝试中，只需要一个成功，
其余尝试的最佳策略是尽可能快速且
无痛地失败。
否则，在大型系统上锁竞争会变得
病态（这里的``大型''是数百个CPU）。
因此，C++17~\cite{RichardSmith2019N4800}在其\co{unique_lock}类中
提供了从严格RAII中的逃逸，允许critical section的范围
被控制到与使用显式锁获取和释放原语
大致相同的程度。

\begin{figure}
\centering
\resizebox{3in}{!}{\includegraphics{locking/rnplock}}
\caption{加锁层级}
\label{fig:locking:Locking Hierarchy}
\end{figure}

严格RAII不友好的数据结构的示例来自Linux内核的RCU，
如\cref{fig:locking:Locking Hierarchy}所示。
每个CPU被分配一个叶\co{rcu_node}结构，每个\co{rcu_node}结构
都有一个指向其父节点的指针（名称颇为直白，叫\co{->parent}），
一直向上到根\co{rcu_node}结构，其\co{->parent}指针为\co{NULL}。
每个父节点拥有的子\co{rcu_node}结构数量可以不同，
但通常为32或64个。
每个\co{rcu_node}结构还包含一个名为\co{->fqslock}的锁。

总体方法是一个\emph{锦标赛}，
给定的CPU有条件地获取其
叶子\co{rcu_node}结构的\co{->fqslock}，如果成功，
则尝试获取父节点的，然后释放子节点的。
此外，在每一层，CPU检查全局\co{gp_flags}
变量，如果该变量表明某个其他CPU已经
断言了该事件，则第一个CPU退出竞争。
这种先获取后释放的序列一直持续到
\co{gp_flags}变量表明其他人赢得了锦标赛、
获取\co{->fqslock}的某次尝试失败、或者
根\co{rcu_node}结构的\co{->fqslock}已被获取为止。
如果根\co{rcu_node}结构的\co{->fqslock}被获取，
则调用名为\co{do_force_quiescent_state()}的函数。

\begin{listing}
\begin{fcvlabel}[ln:locking:Conditional Locking to Reduce Contention]
\begin{VerbatimL}[commandchars=\\\[\]]
void force_quiescent_state(struct rcu_node *rnp_leaf)
{
	int ret;
	struct rcu_node *rnp = rnp_leaf;
	struct rcu_node *rnp_old = NULL;

	for (; rnp != NULL; rnp = rnp->parent) {	\lnlbl[loop:b]
		ret = (READ_ONCE(gp_flags)) ||		\lnlbl[flag_set]
		       !raw_spin_trylock(&rnp->fqslock);\lnlbl[trylock]
		if (rnp_old != NULL)			\lnlbl[non_NULL]
			raw_spin_unlock(&rnp_old->fqslock);\lnlbl[rel1]
		if (ret)				\lnlbl[giveup]
			return;				\lnlbl[return]
		rnp_old = rnp;				\lnlbl[save]
	}						\lnlbl[loop:e]
	if (!READ_ONCE(gp_flags)) {			\lnlbl[flag_not_set]
		WRITE_ONCE(gp_flags, 1);		\lnlbl[set_flag]
		do_force_quiescent_state();		\lnlbl[invoke]
		WRITE_ONCE(gp_flags, 0);		\lnlbl[clr_flag]
	}
	raw_spin_unlock(&rnp_old->fqslock);		\lnlbl[rel2]
}
\end{VerbatimL}
\end{fcvlabel}
\caption{通过条件加锁减少竞争}
\label{lst:locking:Conditional Locking to Reduce Contention}
\end{listing}

实现此功能的简化代码如
\cref{lst:locking:Conditional Locking to Reduce Contention}所示。
这个函数的目的是在并发检测到
需要调用\co{do_force_quiescent_state()}函数的CPU之间进行调解。
在任何给定时间，只有一个\co{do_force_quiescent_state()}
实例处于活动状态才有意义，因此如果有多个
并发调用者，我们最多需要其中一个实际调用
\co{do_force_quiescent_state()}，而其余的需要（尽可能快速且
无痛地）放弃并离开。

\begin{fcvref}[ln:locking:Conditional Locking to Reduce Contention]
为此，循环\clnrefrange{loop:b}{loop:e}的每次遍历尝试
在\co{rcu_node}层级中向上推进一层。
如果\co{gp_flags}变量已经设置（\clnref{flag_set}）或者尝试
获取当前\co{rcu_node}结构的\co{->fqslock}
不成功（\clnref{trylock}），则局部变量\co{ret}被设为1。
如果\clnref{non_NULL}看到局部变量\co{rnp_old}非\co{NULL}，
意味着我们持有\co{rnp_old}的\co{->fqs_lock}，
\clnref{rel1}释放此lock（但仅在尝试获取
父\co{rcu_node}结构的\co{->fqslock}之后）。
如果\clnref{giveup}看到\clnref{flag_set}或~\lnref{trylock}
有放弃的理由，
\clnref{return}返回给调用者。
否则，我们必定已获取当前\co{rcu_node}结构的
\co{->fqslock}，所以\clnref{save}将指向此结构的指针保存在
局部变量\co{rnp_old}中，为循环的下一次遍历做准备。

如果控制到达\clnref{flag_not_set}，我们赢得了锦标赛，现在持有
根\co{rcu_node}结构的\co{->fqslock}。
如果\clnref{flag_not_set}仍然看到全局变量\co{gp_flags}为零，
\clnref{set_flag}将\co{gp_flags}设为1，\clnref{invoke}调用
\co{do_force_quiescent_state()}，
\clnref{clr_flag}将\co{gp_flags}重置为零。
无论如何，\clnref{rel2}释放根\co{rcu_node}结构的
\co{->fqslock}。
\end{fcvref}

\QuickQuizSeries{%
\QuickQuizB{
	\cref{lst:locking:Conditional Locking to Reduce Contention}
	中的代码简直荒谬地复杂！
	为什么不有条件地获取单个全局锁？
}\QuickQuizAnswerB{
	有条件地获取单个全局lock确实效果很好，
	但仅适用于相对较少的CPU数量。
	要了解为什么在拥有数百个CPU的系统中
	这是有问题的，请看
	\cref{fig:count:Atomic Increment Scalability on x86}。
}\QuickQuizEndB
%
\QuickQuizE{
	\begin{fcvref}[ln:locking:Conditional Locking to Reduce Contention]
	等一下！
	如果我们在\cref{lst:locking:Conditional Locking to Reduce Contention}
	的\clnref{flag_not_set}上``赢得''了锦标赛，
	我们得做\co{do_force_quiescent_state()}的所有工作。
	这到底算什么赢，真的？
	\end{fcvref}
}\QuickQuizAnswerE{
	确实如此？
	这只是表明在并发中，就像在生活中一样，
	在参加比赛之前应该注意了解
	赢意味着什么。
}\QuickQuizEndE
}

这个函数展示了并不罕见的分层加锁模式。
这种模式使用严格的RAII加锁难以实现，\footnote{
	这就是为什么许多RAII加锁实现提供了一种方式
	将lock从获取它的作用域泄漏到
	将要释放它的作用域中。
	然而，必须有某个对象来调解作用域泄漏，这与
	非RAII的显式加锁原语相比可能增加复杂性。}
就像前面提到的迭代器封装一样，因此显式的
锁/解锁原语（或C++17风格的\co{unique_lock}逃逸）将
在可预见的未来被需要。

\section{加锁实现问题}
\label{sec:locking:Locking Implementation Issues}
%
\epigraph{当你把梦想转化为现实时，它永远不是一个完整的实现。
	  做梦比做事容易。}
	 {Shai Agassi}

开发者几乎总是最好使用系统提供的任何加锁原语，例如POSIX
pthread mutex locks~\cite{OpenGroup1997pthreads,Butenhof1997pthreads}.
尽管如此，研究示例实现可能会有帮助，
考虑极端工作负载和环境带来的挑战也是如此。

\subsection{基于原子交换的互斥锁示例实现}
\label{sec:locking:Sample Exclusive-Locking Implementation Based on Atomic Exchange}

\begin{fcvref}[ln:locking:xchglock:lock_unlock]
本节回顾\cref{lst:locking:Sample Lock Based on Atomic Exchange}中
展示的实现。
这个锁的数据结构只是一个\co{int}，如
\clnref{typedef}所示，但可以是任何整数类型。
这个锁的初始值为零，意思是``未锁定''，
如\clnref{initval}所示。
\end{fcvref}

\begin{listing}
\input{CodeSamples/locking/xchglock@lock_unlock.fcv}
\caption{基于原子交换的Lock示例}
\label{lst:locking:Sample Lock Based on Atomic Exchange}
\end{listing}

\QuickQuiz{
	\begin{fcvref}[ln:locking:xchglock:lock_unlock]
	为什么不依赖C语言的默认零初始化，
	而是使用\cref{lst:locking:Sample Lock Based on Atomic Exchange}
	\clnref{initval}上显示的显式初始化器？
	\end{fcvref}
}\QuickQuizAnswer{
	因为这个默认初始化不适用于
	在函数作用域内分配为auto变量的锁。
}\QuickQuizEnd

\begin{fcvref}[ln:locking:xchglock:lock_unlock:lock]
锁获取由\co{xchg_lock()}函数执行，
如\clnrefrange{b}{e}所示。
该函数使用嵌套循环，外层循环反复
将锁的值与值1（意思是``已锁定''）
进行原子交换。
如果旧值已经是1（换句话说，其他人
已经持有该锁），则内层循环（\clnrefrange{inner:b}{inner:e}）
自旋等待直到锁可用，此时外层循环
再次尝试获取锁。
\end{fcvref}

\QuickQuiz{
	\begin{fcvref}[ln:locking:xchglock:lock_unlock:lock]
	为什么要费心使用\cref{lst:locking:Sample Lock Based on Atomic Exchange}
	\clnrefrange{inner:b}{inner:e}上的内层循环？
	为什么不简单地在\clnref{atmxchg}上
	反复执行原子交换操作？
	\end{fcvref}
}\QuickQuizAnswer{
	\begin{fcvref}[ln:locking:xchglock:lock_unlock:lock]
	假设锁被持有，并且有几个线程
	正在尝试获取该锁。
	在这种情况下，如果这些线程都在原子交换操作上
	循环，它们将在它们之间来回弹跳包含锁的
	缓存行，给互联施加负载。
	相比之下，如果这些线程在
	\clnrefrange{inner:b}{inner:e}的内层循环中自旋，
	它们将各自在自己的缓存中自旋，
	给互联施加的负载可以忽略不计。
	\end{fcvref}
}\QuickQuizEnd

\begin{fcvref}[ln:locking:xchglock:lock_unlock:unlock]
锁释放由\co{xchg_unlock()}函数执行，
如\clnrefrange{b}{e}所示。
\Clnref{atmxchg}将值零（``未锁定''）原子地交换到
锁中，从而将其标记为已释放。
\end{fcvref}

\QuickQuiz{
	\begin{fcvref}[ln:locking:xchglock:lock_unlock:unlock]
	为什么不简单地在\cref{lst:locking:Sample Lock Based on Atomic Exchange}
	的\clnref{atmxchg}上将零存储到锁字中？
	\end{fcvref}
}\QuickQuizAnswer{
	这可以是一个合法的实现，但前提是
	这个存储之前有一个memory barrier并使用
	了\co{WRITE_ONCE()}。
	当使用\co{xchg()}操作时不需要memory barrier，
	因为该操作由于返回一个值而隐含了
	一个完整的memory barrier。
}\QuickQuizEnd

这个锁是test-and-set lock~\cite{Segall84}的一个简单例子，
但非常类似的机制已在生产中被广泛
用作纯spinlock。

\subsection{其他互斥锁实现}
\label{sec:locking:Other Exclusive-Locking Implementations}

基于原子指令还有很多其他可能的加锁实现，
许多在Mellor-Crummey和Scott的经典论文~\cite{MellorCrummey91a}中有回顾。
这些实现代表了多维设计权衡中的不同
点~\cite{Guerraoui:2019:LPA:3319851.3301501,HugoGuirouxPhD,McKenney96a}。
例如，前一节介绍的基于原子交换的test-and-set锁
在竞争较低时工作良好，且具有内存占用小的优点。
它避免将锁授予无法使用它的线程，但因此
在高竞争级别下可能遭受\IX{unfairness}甚至\IX{starvation}。

相比之下，ticket锁~\cite{MellorCrummey91a}曾在Linux内核中使用，
避免了高竞争级别下的unfairness。
然而，由于其严格的FIFO原则，它可以将锁授予
当前无法使用它的线程，也许是因为该线程被抢占或中断。
另一方面，重要的是不要对抢占和中断的可能性过于担心。
毕竟，在许多情况下，这种抢占和中断同样可能发生在
锁被获取之后。\footnote{
	此外，处理高锁竞争的最好方法是首先避免它！
	然而，有些情况下高锁竞争是可用的弊端中
	最小的一个，而且无论如何，研究处理高竞争级别的方案
	是一个很好的思维练习。}

所有等待者在单个内存位置上自旋的加锁实现，
包括test-and-set锁和ticket锁，
在高竞争级别下都存在性能问题。
问题是释放锁的线程必须更新相应内存位置的值。
在低竞争下，这不是问题：
相应的缓存行很可能仍然是持有锁的线程的本地且可写。
相比之下，在高竞争级别下，每个试图获取锁的线程
都将有\IX{cache line}的只读副本，而锁持有者
需要在执行释放锁的更新之前使所有这些副本无效。
一般来说，CPU和线程越多，在高竞争条件下
释放锁时产生的开销就越大。

这种负向可扩展性推动了许多不同的排队锁
实现~\cite{Anderson90,Graunke90,MellorCrummey91a,Wisniewski94,Craig93,Magnusson94,Takada93}，
其中一些在最近版本的Linux内核中使用~\cite{JonathanCorbet2014qspinlocks}。
排队锁通过为每个线程分配一个队列元素来避免高缓存无效化开销。
这些队列元素被链接在一起形成一个队列，该队列管理着
锁将按照什么顺序被授予等待线程。
关键点是每个线程在自己的队列元素上自旋，
因此锁持有者只需要使下一个线程CPU缓存中的
第一个元素无效。
这种安排大大降低了高竞争级别下锁移交的开销。

较新的排队锁实现还会考虑系统架构，优先在本地授予锁，
同时采取措施避免
饥饿~\cite{McKenney02e,radovic03hierarchical,radovic02efficient,BenJackson02,McKenney02d}。
其中许多可以类比为传统磁盘I/O调度中使用的电梯算法。
Dice等人~讨论了使用本地和全局锁来将任意不感知架构的锁
转换为本地/全局加锁方案，通过提供（例如）每socket的锁
加上全局锁来优化系统
结构~\cite{DavidDice2012NUMAlock1,DavidDice2012NUMAlock2}。

不幸的是，同样的调度逻辑虽然在高竞争时提高了排队锁的
\IX{efficiency}，但在低竞争时却增加了其开销。
Beng-Hong Lim和Anant Agarwal因此将简单的test-and-set锁与
排队锁结合起来，在低竞争时使用test-and-set锁，
在高竞争时切换到排队锁~\cite{BengHongLim94}，
从而在低竞争时获得低开销，在高竞争时获得公平性和高吞吐量。
Browning等人~采取了类似的方法，但避免使用单独的标志，
使得test-and-set快速路径使用与简单test-and-set锁
相同的指令序列~\cite{LukeBrowning2005SimpleLockNUMAAware}。
这种方法已在生产环境中使用。

在高竞争级别下出现的另一个问题是锁持有者被延迟，
特别是当延迟是由于抢占引起时，这可能导致\emph{优先级反转}，
即低优先级线程持有锁，但被
中等优先级的CPU密集型线程抢占，导致
高优先级进程在尝试获取锁时阻塞。
结果是CPU密集型的中等优先级进程阻止了
高优先级进程运行。
一种解决方案是\emph{优先级继承}~\cite{Lampson1980Mesa}，
它已被广泛用于实时
计算~\cite{LuiSha1990PriorityInheritance,JonathanCorbet2006PriorityInheritance}，
尽管对这种做法仍有一些争议~\cite{Yodaiken2004FSM,DougLocke2002a}。

避免优先级反转的另一种方法是在持有锁时
防止抢占。
因为在持有锁时防止抢占也提高了吞吐量，
大多数专有UNIX内核提供了某种形式的调度器感知
同步机制~\cite{Kontothanassis97a}，
这在很大程度上归功于某个相当大的数据库供应商的努力。
这些机制通常采用在给定代码区域中
应避免抢占的提示形式，该提示通常
被放在机器寄存器中。
这些提示通常采用在特定机器寄存器中设置位的形式，
这使得这些机制具有极低的每次锁获取开销。
相比之下，Linux避免这些提示。
取而代之的是，Linux内核社区对调度器感知同步请求的
回应是一种称为\emph{futex}的机制~\cite{HubertusFrancke2002Futex,IngoMolnar2006RobustFutexes,StevenRostedt2006piFutexes,UlrichDrepper2011Futexes}。

有趣的是，实现锁并不严格需要原子指令~\cite{Dijkstra65a,Lamport74a}。
关于基于简单load和store的加锁实现相关问题的
优秀阐述可以在Herlihy和Shavit的教科书~\cite{HerlihyShavit2008Textbook,HerlihyShavit2020Textbook}中找到。
这里重复的要点是，此类实现目前
几乎没有实际应用，尽管仔细研究
它们既有趣又有启发性。
然而，除了下面描述的一个例外，此类研究留给
读者作为练习。

Gamsa等人~\cite[Section 5.3]{Gamsa99}描述了一种基于令牌的
机制，其中令牌在CPU之间循环。
当令牌到达给定CPU时，该CPU对
由该令牌保护的任何东西具有排他访问权。
有许多方案可用于实现基于令牌的机制，
例如：

\begin{enumerate}
\item	维护一个每CPU标志，初始时除一个CPU外其余均为零。
	当CPU的标志非零时，它持有令牌。
	当它用完令牌后，将自己的标志清零，
	并将下一个CPU的标志设为一（或任何其他非零值）。
\item	维护一个每CPU计数器，初始时设置为
	对应CPU的编号，我们假设编号范围为零到$N-1$，
	其中$N$是系统中的CPU数量。
	当一个CPU的计数器大于下一个CPU的计数器
	（考虑计数器回绕）时，该CPU持有令牌。
	当它用完令牌后，将下一个CPU的计数器
	设为比自己计数器大一的值。
\end{enumerate}

\QuickQuizSeries{%
\QuickQuizB{
	如何判断一个计数器是否大于另一个，同时考虑计数器回绕？
}\QuickQuizAnswerB{
	在C语言中，以下宏可以正确处理这一问题：

\begin{VerbatimU}
#define ULONG_CMP_LT(a, b) \
        (ULONG_MAX / 2 < (a) - (b))
\end{VerbatimU}

	虽然直接对两个有符号整数做减法很诱人，
	但应当避免这样做，因为在C语言中，有符号整数溢出是未定义行为。
	例如，如果编译器知道其中一个值为正，另一个为负，
	它完全有权直接假设正数大于负数，
	即使正数减去负数很可能导致溢出从而得到负数。

	编译器如何能知道两个数的符号呢？
	它可能根据之前的赋值和比较推断出来。
	在这种情况下，如果每CPU计数器是有符号的，
	编译器可以推断它们的值一直在增加，
	然后可能假设它们永远不会变为负数。
	这种假设很可能导致编译器生成
	有问题的代码~\cite{PaulEMcKenney2012SignedOverflow,JohnRegehr2010UndefinedBehavior}。
}\QuickQuizEndB
%
\QuickQuizE{
	哪种方式更好，计数器方法还是标志方法？
}\QuickQuizAnswerE{
	标志方法通常会产生更少的缓存未命中，
	但更好的答案是两种都试一试，
	看看哪种对你的特定工作负载效果最好。
}\QuickQuizEndE
}

这种锁的不寻常之处在于给定CPU不一定能立即获取它，
即使当前没有其他CPU在使用它。
相反，CPU必须等待令牌轮到它。
这在CPU需要定期访问\IX{critical
section}但能容忍令牌循环速率变化的情况下很有用。
Gamsa等人~\cite{Gamsa99}用它来实现read-copy update的一个变体
（参见\cref{sec:defer:Read-Copy Update (RCU)}），
但它也可以用于保护周期性的每CPU操作，例如
刷新内存分配器使用的每CPU缓存~\cite{McKenney93}、
垃圾回收每CPU数据结构，或将每CPU数据
刷新到共享存储（或大容量存储）。

Linux内核现在使用排队spinlock~\cite{JonathanCorbet2014qspinlocks}，
但由于在竞争级别范围内提供良好性能的实现的复杂性，
这条道路并不总是一帆风顺~\cite{CatalinMarinas2018qspinlockTLA,WillDeacon2018qspinlock}。
随着越来越多的人熟悉并行硬件
并并行化越来越多的代码，我们可以继续期待更多
特殊用途的加锁原语出现，例如参见Guerraoui等人~\cite{Guerraoui:2019:LPA:3319851.3301501,HugoGuirouxPhD}。
尽管如此，你应该仔细考虑这个重要的安全提示：
尽可能使用标准同步原语。
标准同步原语相对于自己动手的努力的
巨大优势在于标准原语通常
\emph{大大}减少了bug倾向。\footnote{
	是的，我至少做了自己应有份额的自制
	同步原语。
	然而，你会注意到我的头发比
	我开始做那种工作之前灰白得多。
	巧合？
	也许吧。
	但你\emph{真的}愿意冒着自己的头发
	过早变灰的风险吗？}

% locking/locking-existence.tex
% mainfile: ../perfbook.tex
% SPDX-License-Identifier: CC-BY-SA-3.0

\section{基于锁的存在性保证}
\label{sec:locking:Lock-Based Existence Guarantees}
%
\epigraph{存在先于本质并支配本质。}{Jean-Paul Sartre}

并行编程中的一个关键挑战是提供
\emph{\IXpl{existence guarantee}}~\cite{Gamsa99}，
以便对给定对象的访问尝试可以依赖于该对象在整个访问尝试期间一直存在。

\begin{listing}
\begin{fcvlabel}[ln:locking:Per-Element Locking Without Existence Guarantees]
\begin{VerbatimL}[commandchars=\\\@\$]
int delete(int key)
{
	int b;
	struct element *p;

	b = hashfunction(key);
	p = hashtable[b];
	if (p == NULL || p->key != key)		\lnlbl@chk_first$
		return 0;
	spin_lock(&p->lock);			\lnlbl@acq$
	hashtable[b] = NULL;			\lnlbl@NULL$
	spin_unlock(&p->lock);
	kfree(p);
	return 1;				\lnlbl@return1$
}
\end{VerbatimL}
\end{fcvlabel}
\caption{无存在性保证的逐元素加锁（有缺陷！）}
\label{lst:locking:Per-Element Locking Without Existence Guarantees (Buggy!)}
\end{listing}

在某些情况下，存在性保证是隐式的：

\begin{enumerate}
\item	基础模块中的全局变量和静态局部变量将在应用程序运行期间一直存在。
\item	已加载模块中的全局变量和静态局部变量将在该模块保持加载状态期间一直存在。
\item	只要模块的某个函数有活跃实例，该模块就会保持加载状态。
\item	给定函数实例的栈上变量将一直存在，直到该实例返回。
\item	如果你正在某个函数中执行，或者被该函数直接或间接调用，
	那么该函数就有一个活跃实例。
\end{enumerate}

这些隐式存在性保证很直观，但涉及隐式存在性保证的
bug 确实可能发生。

\QuickQuiz{
	依赖隐式存在性保证怎么会导致 bug？
}\QuickQuizAnswer{
	以下是一些因不当使用隐式存在性保证而导致的 bug：
	\begin{enumerate}
	\item	程序将一个全局变量的地址写入文件，
		然后该程序的后续实例读取该地址并尝试解引用它。
		由于地址空间随机化，这可能会失败，
		更不用说程序的重新编译了。
	\item	一个模块可以将其某个变量的地址记录在
		位于其他模块中的指针中，然后在该模块被
		卸载后尝试解引用该指针。
	\item	一个函数可以将其某个栈上变量的地址记录到
		一个全局指针中，其他函数可能在该函数
		返回后尝试解引用该指针。
	\end{enumerate}
	我相信你能想出更多的可能性。
}\QuickQuizEnd

但更有趣——也更棘手——的保证涉及堆内存：
动态分配的数据结构将一直存在，直到它被释放。
要解决的问题是如何将该结构的释放与对同一结构的并发访问同步。
一种方法是使用\emph{显式保证}，例如加锁。
如果只有在持有某个锁时才能释放给定结构，那么持有该锁就保证了该结构的存在。

但这个保证依赖于锁本身的存在。
保证锁存在的一种直接方法是将锁放在全局变量中，但全局加锁的缺点是限制了可扩展性。
一种随数据结构大小增长而提高可扩展性的方法是在结构的每个元素中放置一个锁。
不幸的是，将用于保护数据元素的锁放在数据元素本身中
容易产生微妙的\IXpl{race condition}，如
\cref{lst:locking:Per-Element Locking Without Existence Guarantees (Buggy!)}
所示。

\QuickQuiz{
	\begin{fcvref}[ln:locking:Per-Element Locking Without Existence Guarantees]
	如果我们需要删除的元素不是
	\cref{lst:locking:Per-Element Locking Without Existence Guarantees (Buggy!)}
	\clnref{chk_first} 上列表的第一个元素怎么办？
	\end{fcvref}
}\QuickQuizAnswer{
	这是一个没有链接的非常简单的哈希表，所以给定桶中唯一的
	元素就是第一个元素。
	读者可以尝试将此示例改编为具有完整链接的哈希表。
}\QuickQuizEnd


\begin{fcvref}[ln:locking:Per-Element Locking Without Existence Guarantees]
要看到其中一个 race condition，请考虑以下事件序列：
\begin{enumerate}
\item	线程~0 调用 \co{delete(0)}，执行到代码清单的 \clnref{acq}，获取了锁。
\item	线程~1 并发调用 \co{delete(0)}，也执行到 \clnref{acq}，
	但因线程~0 持有锁而自旋等待。
\item	线程~0 执行 \clnrefrange{NULL}{return1}，将元素从
	哈希表中移除，释放锁，然后释放该元素。
\item	线程~0 继续执行，并分配内存，恰好得到了
	它刚刚释放的那块内存。
\item	线程~0 然后将这块内存初始化为
	其他类型的结构。
\item	线程~1 的 \co{spin_lock()} 操作失败，因为
	它认为是 \co{p->lock} 的东西已经不再
	是一个 spinlock 了。
\end{enumerate}
因为没有存在性保证，当线程试图在 \clnref{acq} 上获取
该元素的锁时，数据元素的身份可能已经改变了！
\end{fcvref}

\begin{listing}
\begin{fcvlabel}[ln:locking:Per-Element Locking With Lock-Based Existence Guarantees]
\begin{VerbatimL}[commandchars=\\\@\$]
int delete(int key)
{
	int b;
	struct element *p;
	spinlock_t *sp;

	b = hashfunction(key);
	sp = &locktable[b];
	spin_lock(sp);				\lnlbl@acq$
	p = hashtable[b];			\lnlbl@getp$
	if (p == NULL || p->key != key) {
		spin_unlock(sp);
		return 0;
	}
	hashtable[b] = NULL;
	spin_unlock(sp);
	kfree(p);
	return 1;
}
\end{VerbatimL}
\end{fcvlabel}
\caption{具有基于锁的存在性保证的逐元素加锁}
\label{lst:locking:Per-Element Locking With Lock-Based Existence Guarantees}
\end{listing}

\begin{fcvref}[ln:locking:Per-Element Locking With Lock-Based Existence Guarantees]
修复此示例的一种方法是使用一组哈希化的全局锁，使每个哈希桶都有自己的锁，如
\cref{lst:locking:Per-Element Locking With Lock-Based Existence Guarantees}
所示。
这种方法允许在获取数据元素的指针之前（在 \clnref{getp}）先获取适当的锁（在 \clnref{acq}）。
虽然这种方法对于包含在单个可分区数据结构（如代码清单所示的哈希表）中的元素效果很好，
但如果给定数据元素可以是多个哈希表的成员，
或者对于更复杂的数据结构如树或图，这种方法可能会有问题。
这些问题不仅可以解决，而且其解决方案还构成了基于锁的软件事务内存
实现的基础~\cite{Shavit95,DaveDice2006DISC}。
然而，
\cref{chp:Deferred Processing}
描述了提供存在性保证的更简单——也更快——的方法。
\end{fcvref}

% @@@ Optimized sharded locking from P1726R5, and pointer-zap implications.


\section{加锁：
		  英雄还是反派？}
\label{sec:locking:Locking: Hero or Villain?}
%
\epigraph{你要么作为英雄而死，要么活得够久，看到自己变成反派。}
	 {Aaron Eckhart \emph{饰演} Harvey Dent}

正如现实生活中经常发生的那样，加锁可以是英雄也可以是反派，
取决于它的使用方式和手头的问题。
根据我的经验，编写整个应用程序的人对加锁感到满意，
编写并行库的人不太满意，而将现有顺序库并行化的人
则极其不满意。
以下各节讨论这些观点差异的一些原因。

\subsection{应用程序的加锁：
				      英雄！}
\label{sec:locking:Locking For Applications: Hero!}

当编写整个应用程序（或整个内核）时，开发者对设计拥有完全控制，
包括同步设计。
假设设计充分利用了分区，如
\cref{chp:Partitioning and Synchronization Design}中所讨论的，
加锁可以是一种极其有效的同步机制，
这从生产级并行软件中大量使用加锁中得到了证明。

然而，尽管这类软件通常将大部分同步设计基于加锁，
它也几乎总是使用其他同步机制，包括
特殊计数算法（\cref{chp:Counting}）、
数据所有权（\cref{chp:Data Ownership}）、
引用计数（\cref{sec:defer:Reference Counting}）、
hazard pointers（\cref{sec:defer:Hazard Pointers}）、
sequence locking（\cref{sec:defer:Sequence Locks}）和
read-copy update（\cref{sec:defer:Read-Copy Update (RCU)}）。
此外，从业者使用死锁
检测~\cite{JonathanCorbet2006lockdep}、
锁获取/释放平衡~\cite{JonathanCorbet2004sparse}、
缓存未命中分析~\cite{ValgrindHomePage}、
基于硬件计数器的性能分析~\cite{LinuxKernelPerfWiki,OProfileHomePage}
等工具以及许多其他工具。

有了仔细的设计、使用好的同步机制组合和好的工具，
加锁对应用程序和内核效果很好。

\subsection{并行库的加锁：
					    只是另一个工具}
\label{sec:locking:Locking For Parallel Libraries: Just Another Tool}

与应用程序和内核不同，库的设计者无法知道库将与之交互的代码的
加锁设计。
事实上，那些代码可能多年后才会编写。
因此，库设计者有更少的控制权，在制定其同步设计时
必须更加小心。

死锁当然是特别令人关注的，\cref{sec:locking:Deadlock}中
讨论的技术需要被应用。
因此，一种流行的死锁避免策略是确保库的锁是
封装程序的加锁层次结构的独立子树。
然而，这可能比看起来更难。

一个复杂情况在\cref{sec:locking:Local Locking Hierarchies}中讨论过，
即当库函数调用应用程序代码时，\co{qsort()}的
比较函数参数就是一个典型的例子。
另一个复杂情况是与信号处理程序的交互。
如果应用程序的信号处理程序由在库函数中接收到的信号调用，
那么死锁就会像库函数直接调用信号处理程序一样确定地发生。
最后一个复杂情况出现在那些可以在\co{fork()}/\co{exec()}对之间
使用的库函数中，例如由于使用\co{system()}函数。
在这种情况下，如果你的库函数在\co{fork()}时持有锁，
那么子进程将在该锁被持有的情况下开始生命。
由于释放该锁的线程运行在父进程中而不在子进程中，
如果子进程调用你的库函数，死锁就会发生。

以下策略可用于避免这些情况下的死锁问题：

\begin{enumerate}
\item	不使用回调也不使用信号。
\item	不在回调或信号处理程序中获取锁。
\item	让调用者控制同步。
\item	参数化库API，将加锁委托给调用者。
\item	显式避免回调死锁。
\item	显式避免信号处理程序死锁。
\item	避免调用\co{fork()}。
\end{enumerate}

以上每种策略将在以下各节中讨论。

\subsubsection{既不使用回调也不使用信号}
\label{sec:locking:Use Neither Callbacks Nor Signals}

如果一个库函数避免使用回调，且整个应用程序避免使用信号，
那么该库函数获取的任何锁都将是加锁层级树的叶节点。
这种安排可以避免死锁，如
\cref{sec:locking:Locking Hierarchies}中所讨论的。
虽然该策略在适用的地方效果极好，
但有些应用程序必须使用信号处理程序，
有些库函数（如
\cref{sec:locking:Local Locking Hierarchies}中讨论的\co{qsort()}函数）
需要使用回调。

下一节描述的策略通常可以用于这些情况。

\subsubsection{避免在回调和信号处理程序中加锁}
\label{sec:locking:Avoid Locking in Callbacks and Signal Handlers}

如果回调和信号处理程序都不获取锁，
那么它们就不会被卷入死锁循环中，
这使得简单的加锁层级可以再次将库函数视为
加锁层级树的叶节点。
该策略对\co{qsort}的大多数使用场景效果极好，
因为其回调通常只是简单地比较传入的两个值。
该策略对许多信号处理程序也同样出色，
特别是考虑到从信号处理程序中获取锁
通常是不推荐的~\cite{TheOpenGroup1997SUS},\footnote{
	但标准的规定并不能阻止聪明的程序员使用原子操作
	自造加锁原语。}
但如果应用程序需要在信号处理程序中操作复杂数据结构，
该策略可能失效。

以下是在必须操作复杂数据结构时避免在信号处理程序中获取锁的一些方法：

\begin{enumerate}
\item	使用基于\IXacrl{nbs}的简单数据结构，
	如
	\cref{sec:advsync:Simple NBS}中将讨论的。
\item	如果数据结构过于复杂，不适合合理使用
	非阻塞同步，则创建一个允许
	非阻塞入队操作的队列。
	在信号处理程序中，不操作复杂数据结构，
	而是向队列中添加一个描述所需更改的元素。
	然后一个单独的线程可以从队列中取出元素，
	使用普通的加锁方式执行所需的更改。
	有许多现成的并发队列
	实现~\cite{ChristophMKirsch2012FIFOisntTR,MathieuDesnoyers2009URCU,MichaelScott96}。
\end{enumerate}

该策略应通过偶尔的手动或（更好的）自动检查
回调和信号处理程序来执行。
在进行这些检查时，要警惕那些可能
（不明智地）使用原子操作自造锁的聪明程序员。

\subsubsection{调用者控制同步}
\label{sec:locking:Caller Controls Synchronization}

让调用者控制同步在库函数操作数据结构的
独立调用者可见实例时效果极好，
每个实例都可以单独同步。
例如，如果库函数操作搜索树，
而应用程序需要大量独立的搜索树，
那么应用程序可以为每棵树关联一个锁。
然后应用程序根据需要获取和释放锁，
库根本无需感知并行性。
由应用程序控制并行性，
因此加锁可以工作得很好，
如\cref{sec:locking:Locking For Applications: Hero!}中所讨论的。

然而，如果库实现了一个需要内部并发的数据结构，
例如哈希表或并行排序，
该策略就会失效。
在这种情况下，库必须控制自己的同步。

\subsubsection{参数化库同步}
\label{sec:locking:Parameterize Library Synchronization}

其思想是在库的API中添加参数来指定
获取哪些锁、如何获取和释放它们，或两者兼有。
该策略允许应用程序承担避免死锁的全局任务，
通过指定获取哪些锁（传入相关锁的指针）
以及如何获取它们（传入锁获取和释放函数的指针），
同时也允许给定的库函数通过决定在哪里获取和释放锁
来控制自己的并发性。

特别地，该策略允许锁获取和释放函数根据需要屏蔽信号，
而库代码无需关心哪些信号需要被哪些锁屏蔽。
该策略使用的关注点分离可以非常有效，
但在某些情况下，以下各节中提出的策略可能效果更好。

话虽如此，向外部API传递显式锁指针必须非常谨慎，
如\cref{sec:locking:Locking Hierarchies and Pointers to Locks}中所讨论的。
虽然这种做法有时是正确的，
但你应该先研究替代设计来帮自己一个忙。

\subsubsection{显式避免回调死锁}
\label{sec:locking:Explicitly Avoid Callback Deadlocks}

该策略背后的基本规则在
\cref{sec:locking:Local Locking Hierarchies}中已讨论过：
``在调用未知代码之前释放所有锁。''
这通常是最佳方法，因为它允许应用程序
忽略库的加锁层级：
库仍然是应用程序整体加锁层级的叶节点或孤立子树。

在无法在调用未知代码之前释放所有锁的情况下，
\cref{sec:locking:Layered Locking Hierarchies}中描述的分层加锁层级可以很好地工作。
例如，如果未知代码是信号处理程序，
这意味着库函数需要在所有锁获取期间屏蔽信号，
这可能既复杂又缓慢。
因此，在信号处理程序（可能不明智地）获取锁的情况下，
下一节中的策略可能有所帮助。

\subsubsection{显式避免信号处理程序死锁}
\label{sec:locking:Explicitly Avoid Signal-Handler Deadlocks}

假设已知某个库函数会获取锁，但不屏蔽信号。
进一步假设需要从信号处理程序内部和外部
都调用该函数，且不允许修改该库函数。
当然，如果不采取特殊措施，
当该库函数持有锁时收到信号，
而信号处理程序又调用了同一个库函数，
该库函数进而尝试重新获取同一把锁，
死锁就会发生。

此类死锁可以通过以下方式避免：

\begin{enumerate}
\item	如果应用程序从信号处理程序内部调用该库函数，
	那么每次从信号处理程序外部调用该库函数时，
	都必须屏蔽该信号。
\item	如果应用程序在持有某个信号处理程序内获取的锁时
	调用该库函数，那么每次在信号处理程序外部
	调用该库函数时，都必须屏蔽该信号。
\end{enumerate}

这些规则可以通过使用类似于
Linux内核lockdep锁依赖检查器的工具来
强制执行~\cite{JonathanCorbet2006lockdep}。
lockdep的一个巨大优势是它不会被
人类直觉所迷惑~\cite{StevenRostedt2011locdepCryptic}。

\subsubsection{在\tco{fork()}和\tco{exec()}之间使用的库函数}
\label{sec:locking:Library Functions Used Between fork() and exec()}

如前所述，如果执行库函数的线程在某个其他线程调用\apipx{fork()}
时持有锁，由于父进程的内存被复制以创建子进程，
该锁在子进程的上下文中将以被持有的状态诞生。
将释放该锁的线程运行在父进程中而不在子进程中，
这意味着虽然父进程的该锁副本会被释放，
但子进程的该锁副本永远不会被释放。
因此，子进程尝试调用同一个库函数
（从而获取同一把锁）的任何企图都会导致死锁。

解决此问题的一种务实且直接的方法是
在进程仍是单线程时\co{fork()}出子进程，
并让该子进程保持单线程。
然后可以将创建更多子进程的请求
传达给这个初始子进程，它就可以安全地
代表其多线程父进程执行任何所需的
\co{fork()}和\apipx{exec()}系统调用。

解决此问题的另一种不那么务实和直接的方法是
让库函数检查锁的所有者是否仍在运行，
如果不在，则通过重新初始化然后获取来``打破''该锁。
然而，这种方法有几个漏洞：

\begin{enumerate}
\item	该锁保护的数据结构可能处于某种中间状态，
	因此简单地打破锁可能导致任意内存损坏。
\item	如果子进程创建了额外的线程，
	两个线程可能同时打破锁，导致
	两个线程都认为自己拥有该锁。
	这同样可能导致任意内存损坏。
\end{enumerate}

\apipx{pthread_atfork()}函数就是为帮助处理这些情况而提供的。
其思想是注册一组三个函数，
一个在\co{fork()}之前由父进程调用，
一个在\co{fork()}之后由父进程调用，
一个在\co{fork()}之后由子进程调用。
然后可以在这三个点执行适当的清理工作。

但请注意，\co{pthread_atfork()}处理程序的编写
通常相当微妙。
\co{pthread_atfork()}效果最好的情况是
相关数据结构可以被子进程简单地重新初始化。
这可能也是POSIX标准禁止在\co{fork()}和\co{exec()}
之间使用任何非异步信号安全函数的原因之一，
这也排除了在此期间获取锁。

\co{fork()}/\co{exec()}的其他替代方案包括\co{posix_spawn()}
和\co{io_uring_spawn()}~\cite{JoshTriplett2022io_uring_spawn,JakeEdge2022io_uring_spawn}。

\subsubsection{并行库：
				   讨论}
\label{sec:locking:Parallel Libraries: Discussion}

无论使用哪种策略，库的API描述必须包含对该策略及调用者
如何与该策略交互的清楚描述。
简而言之，使用加锁构建并行库是可能的，
但不如构建并行应用程序那么容易。

\subsection{顺序库并行化的加锁：
							    反派！}
\label{sec:locking:Locking For Parallelizing Sequential Libraries: Villain!}

随着低成本多核系统的普及，
一项常见任务是将仅为单线程使用设计的现有库并行化。
这种对并行性的忽视太常见了，可能导致从并行编程角度看
严重有缺陷的库API。
候选缺陷包括：

\begin{enumerate}
\item	隐含禁止分区。
\item	需要加锁的回调函数。
\item	面向对象的意大利面条代码。
\end{enumerate}

这些缺陷及其对加锁的影响在以下各节中讨论。

\subsubsection{禁止分区}
\label{sec:locking:Partitioning Prohibited}

假设你正在编写一个单线程哈希表实现。
维护哈希表中元素总数的精确计数既简单又快速，
在每次添加和删除操作时返回这个精确计数也是如此。
那为什么不这样做呢？

原因之一是精确计数器在多核系统上的性能和可扩展性不佳，
如\cref{chp:Counting}中所见。
因此，哈希表的并行化实现将无法
获得良好的性能或可扩展性。

那该怎么办呢？
一种方法是使用\cref{chp:Counting}中的某种算法返回近似计数。
另一种方法是完全放弃元素计数。

无论哪种方式，都需要检查哈希表的使用方式，
以了解添加和删除操作为何需要精确计数。
以下是一些可能的情况：

\begin{enumerate}
\item	确定何时调整哈希表大小。
	在这种情况下，近似计数应该可以很好地工作。
	也可以根据最长链的长度来触发调整大小操作，
	这可以以良好的分区方式按链计算和维护。
\item	估计遍历整个哈希表所需的时间。
	近似计数在这种情况下同样适用。
\item	用于诊断目的，例如检查在向哈希表添加和从中删除元素时
	是否有元素丢失。
	这显然需要精确计数。
	然而，鉴于这种用途本质上是诊断性的，
	可以只维护哈希链的长度，
	然后在屏蔽添加和删除操作的情况下
	不频繁地对它们求和。
\end{enumerate}

事实证明，性能和可扩展性对并行库API的
某些约束现在有了坚实的理论
基础~\cite{HagitAttiya2011LawsOfOrder,Attiya:2011:LOE:1925844.1926442,PaulEMcKenney2011SNC}。
任何设计并行库的人都需要密切关注这些约束。

虽然将真正由于不友好的并发API造成的问题
归咎于加锁是太容易的事，但这并无帮助。
另一方面，人们几乎别无选择，只能同情那些
（比如说）在1985年做出这个选择的不幸开发者。
在那时能预见到并行性需求的开发者是罕见而有勇气的，
而要真正提出一个好的并行友好API，
则需要更加罕见的才华和运气的组合。

时代在变，代码也必须随之改变。
话虽如此，一个流行的库可能有大量的用户，
在这种情况下，对API做出不兼容的更改将是非常愚蠢的。
添加一个并行友好的API来补充现有大量使用的
仅限顺序执行的API通常是最佳做法。

尽管如此，人性如此，我们可以预期
我们的不幸开发者更可能抱怨加锁，
而不是抱怨他自己或她自己糟糕的（虽然可以理解的）API设计选择。

\subsubsection{容易导致死锁的回调}
\label{sec:locking:Deadlock-Prone Callbacks}

\Cref{sec:locking:Local Locking Hierarchies,%
sec:locking:Layered Locking Hierarchies,%
sec:locking:Locking For Parallel Libraries: Just Another Tool}
描述了不受约束地使用回调如何导致加锁问题。
这些章节还描述了如何设计库函数以避免这些问题，
但期望一个没有并行编程经验的1990年代程序员
遵循这样的设计是不现实的。
因此，试图将现有的回调密集型单线程库
并行化的人很可能会有很多机会诅咒加锁的罪恶。

如果回调密集型库有大量使用者，
明智的做法可能是在库中添加并行友好的API，
以允许现有用户逐步转换代码。
或者，有些人提倡在这些情况下使用事务性内存。
虽然事务性内存的前景尚不明朗，
\cref{sec:future:Transactional Memory}讨论了它的优缺点。
需要注意的是，硬件事务性内存
（在\cref{sec:future:Hardware Transactional Memory}中讨论）
在此无法提供帮助，除非硬件事务性内存实现
提供\IXpl{forward-progress guarantee}，而很少有实现能做到。
其他看起来相当实用（即使炒作程度较低）的替代方案包括
\cref{sec:locking:Conditional Locking,sec:locking:Acquire Needed Locks First}
中讨论的方法，以及
\cref{chp:Data Ownership,chp:Deferred Processing}中将讨论的方法。

\subsubsection{面向对象的意大利面条代码}
\label{sec:locking:Object-Oriented Spaghetti Code}

面向对象编程在1980年代或1990年代的某个时候成为主流，
因此生产环境中存在大量的单线程面向对象代码。
虽然面向对象可以是一种有价值的软件技术，
但不受约束地使用对象很容易导致面向对象的意大利面条代码。
在面向对象的意大利面条代码中，控制流以本质上随机的方式在对象之间跳跃，
使代码难以理解，甚至更难（也许不可能）适应加锁层次结构。

虽然许多人可能会认为这样的代码无论如何都应该清理，
但说起来容易做起来难。
如果你被要求将这样的代码并行化，
你可以通过使用\cref{sec:locking:Conditional Locking,sec:locking:Acquire Needed Locks First}
中描述的技术，以及将在\cref{chp:Data Ownership,chp:Deferred Processing}
中讨论的技术来减少咒骂加锁的机会。
这种情况似乎是启发事务内存的用例，
因此也值得尝试一下。
话虽如此，同步机制的选择应该根据\cref{chp:Hardware and its Habits}
中讨论的硬件习惯来决定。
毕竟，如果同步机制的开销比被保护的操作高出几个数量级，
结果不会好看。

这就引出了一个在这些情况下值得提出的问题：
代码应该保持顺序执行吗？
例如，也许应该在进程级别而不是线程级别引入并行性。
一般来说，如果一项任务被证明极其困难，
值得花些时间思考不仅完成该特定任务的替代方法，
而且思考可能更好地解决手头问题的替代任务。

\section{总结}
\label{sec:locking:Summary}
%
\epigraph{成就已解锁。}{佚名}

加锁也许是最广泛使用和最通用的同步工具。
然而，它在从一开始就设计到应用程序或库中时效果最好。
鉴于大量可能需要某天并行运行的已有单线程代码，
加锁因此不应该是你并行编程工具箱中唯一的工具。
接下来几章将讨论其他工具，以及它们如何最好地与加锁
和彼此配合使用。

\QuickQuizAnswersChp{qqzlocking}
